%% 6장 미분방정식과 점화식 (새 목차, 2026-11-25)
\chapter{Differential Equations and Recurrences}
\label{ch:recurrences}

\abstract{특수 함수의 대부분은 이계 선형 미분방정식의 해이고 이웃한 차수끼리 삼항 점화식으로 이어지는데 점화식은 돌리는 방향에 따라 정확하기도 하고 몇 걸음 만에 망가지기도 한다. 이 장은 선형 ODE와 Wronskian, 최소해와 지배해, Miller의 역방향 점화, 연분수, 경계값 문제로 푸는 점화식, Clenshaw 합을 세운다. 구현으로는 Bessel 함수, QUADPACK의 모멘트, Chebyshev 급수를 본다.}


\section{Linear Differential Equations}
\label{sec:ode}

\subsection{Existence and Uniqueness}

구간 $I$에서 연속인 $d\times d$ 행렬 함수 $A(x)$와 벡터 함수 $b(x)$에 대해 선형 연립방정식
\begin{equation}
\label{eq:linear-system}
u'(x)=A(x)u(x)+b(x),\qquad u(x_0)=u_0
\end{equation}
을 생각하자. 이계 방정식 $y''+p(x)y'+q(x)y=0$은 $u=(y,y')^{\mathsf T}$로 두면
\[
u'=\begin{pmatrix}0&1\\-q(x)&-p(x)\end{pmatrix}u
\]
가 되므로 이 꼴에 든다. 선형 방정식에서는 해가 정의역 전체에서 존재한다. 비선형 방정식의 해는 유한한 곳에서 터질 수 있으니 이 성질은 선형성에서 온다.

\begin{theorem}
\label{thm:ode-exist}
$A$, $b$가 구간 $I$에서 연속이면 모든 $x_0\in I$와 $u_0$에 대해 \cref{eq:linear-system}의 해가 $I$ 전체에서 유일하게 존재한다.
\end{theorem}

\begin{proof}
$I$ 안의 닫힌 구간 $J\ni x_0$에서 보이면 충분하다. $J$에서 $\|A(x)\|\le L$이라 하자. \cref{eq:linear-system}은 적분 방정식 $u(x)=u_0+\int_{x_0}^x(A(s)u(s)+b(s))\,ds$와 같다. Picard 반복 $u^{(0)}=u_0$, $u^{(k+1)}(x)=u_0+\int_{x_0}^x(Au^{(k)}+b)\,ds$를 생각하면 차 $\delta^{(k)}=u^{(k+1)}-u^{(k)}$는 $\delta^{(k)}(x)=\int_{x_0}^xA(s)\delta^{(k-1)}(s)\,ds$를 만족한다. $J$에서 $\|\delta^{(0)}\|\le M$이면 귀납적으로
\[
\|\delta^{(k)}(x)\|\le M\frac{(L|x-x_0|)^k}{k!}
\]
이다. 오른쪽의 합이 수렴하므로 $u^{(k)}$는 $J$에서 균등수렴하고 그 극한이 적분 방정식의 연속해이다. 유일성은 두 해의 차 $w$가 $w(x)=\int_{x_0}^xAw\,ds$를 만족한다는 데서 나온다. $\|w\|\le K$로 시작해 같은 부등식을 되풀이하면 $\|w(x)\|\le K(L|x-x_0|)^k/k!\to0$이다.
\end{proof}

이 증명이 비선형 방정식의 경우와 다른 점은 상수 $L$이 해의 크기와 무관하다는 것 하나이고 그 덕분에 반복이 구간 전체에서 한꺼번에 수렴한다. 해 공간의 구조도 이 정리에서 정해진다. $b=0$인 동차 방정식의 해들은 선형공간을 이루고 초깃값 $u(x_0)$로 보내는 사상은 그 공간에서 $\mathbb{F}^d$로 가는 선형 동형이다. 해 공간은 $d$차원이고 이계 방정식은 $2$차원이다.

\subsection{The Wronskian}

해 $d$개를 열로 모은 행렬 $\Phi(x)$를 \term{기본 해 행렬}{fundamental matrix}이라 한다. 해들이 독립인지는 그 행렬식이 $0$인지로 판정한다.

\begin{theorem}[Liouville, Abel]
\label{thm:abel}
동차 방정식 $\Phi'=A\Phi$의 행렬 해에 대해 $\det\Phi(x)=\det\Phi(x_0)\exp\bigl(\int_{x_0}^x\operatorname{tr}A(t)\,dt\bigr)$이다. 특히 이계 방정식 $y''+py'+qy=0$의 두 해의 \term{Wronskian}{Wronskian} $W=y_1y_2'-y_1'y_2$는 $W(x)=W(x_0)\exp(-\int_{x_0}^xp)$이고 $I$의 모든 점에서 $0$이거나 어디서도 $0$이 아니다. 두 해가 일차독립일 필요충분조건은 $W\ne0$이다.
\end{theorem}

\begin{proof}
행렬식의 도함수는 행마다 미분한 행렬식의 합이다. $\Phi'=A\Phi$의 $i$행은 $\sum_ja_{ij}(\Phi\text{의 }j\text{행})$이고 $j\ne i$인 항은 같은 행이 두 번 나오는 행렬식이 되어 사라진다. 그래서 $i$번째 합은 $a_{ii}\det\Phi$이고 $(\det\Phi)'=\operatorname{tr}A\cdot\det\Phi$이다. 이 일계 방정식의 해가 주장한 식이다. 이계 방정식에서는 $\operatorname{tr}A=-p$이고 $\det\Phi=W$이다. 한 점에서 $W=0$이면 두 초깃값 벡터가 종속이고 \cref{thm:ode-exist}의 유일성에서 두 해도 종속이다. 거꾸로 종속이면 $W\equiv0$임은 분명하다.
\end{proof}

\cref{thm:abel} 덕분에 Wronskian은 계산하지 않아도 꼴을 안다. 특수 함수 계산에서는 이 사실을 두 가지로 쓴다. 두 해의 값과 도함수 가운데 셋을 알면 넷째를 Wronskian 관계로 얻을 수 있다. 그리고 계산한 값들이 그 관계를 얼마나 잘 만족하는지로 정확도를 검사할 수 있다. 4--5절에서 stardust의 Bessel 함수는 첫째 용법을 쓴다.

\subsection{Constant Coefficients and the Matrix Exponential}

계수가 상수이면 해는 행렬 지수로 주어진다.

\begin{definition}[행렬 지수]
\label{def:expm}
정사각행렬 $A$에 대해 $e^A=\sum_{k=0}^\infty A^k/k!$로 정의한다.
\end{definition}

급수는 $\|A^k/k!\|\le\|A\|^k/k!$이므로 모든 $A$에서 절대수렴한다. 항별로 미분하면 $(e^{tA})'=Ae^{tA}$이므로 상수 계수 방정식 $u'=Au$, $u(0)=u_0$의 해는 $u(t)=e^{tA}u_0$이고 기본 해 행렬은 $e^{tA}$이다.

\begin{proposition}
\label{prop:expm}
정사각행렬 $A$, $B$에 대해 다음이 성립한다.
\begin{enumerate}
\item $AB=BA$이면 $e^{A+B}=e^Ae^B$이다. 특히 $e^A$는 가역이고 $(e^A)^{-1}=e^{-A}$이다.
\item $A=XJX^{-1}$이면 $e^A=Xe^JX^{-1}$이고 $A$의 고유값이 $\lambda_i$이면 $e^A$의 고유값은 $e^{\lambda_i}$이다.
\item $\det e^A=e^{\operatorname{tr}A}$이다.
\end{enumerate}
\end{proposition}

\begin{proof}
(i) $AB=BA$이면 이항정리 $(A+B)^n=\sum\binom nkA^kB^{n-k}$가 성립하므로 절대수렴하는 두 급수의 곱을 Cauchy 곱으로 정리하면 $e^{A+B}$가 된다. $B=-A$로 두면 둘째 주장이 나온다. (ii) $(XJX^{-1})^k=XJ^kX^{-1}$을 급수에 넣으면 된다. Schur 분해 $A=QTQ^{\mathsf H}$(\cref{sec:schur})에 이를 쓰면 $e^T$는 대각 성분이 $e^{t_{ii}}$인 위삼각 행렬이므로 고유값의 주장이 나온다. (iii) 같은 분해에서 $\det e^A=\prod e^{t_{ii}}=e^{\operatorname{tr}A}$이다. \cref{thm:abel}로도 나온다.
\end{proof}

(i)에서 조건 $AB=BA$를 뺄 수는 없다. 일반적으로 $e^{A+B}\ne e^Ae^B$이며 행렬 지수를 다루는 계산에서 가장 흔한 실수가 여기서 나온다. 행렬 지수는 정의는 간단하지만 계산은 까다롭다. 정의의 급수를 그대로 더하면 $\|A\|$가 클 때 항들이 크게 자랐다가 지워지고 고유값 분해로 계산하면 고유벡터 행렬의 조건수만큼 오차가 커진다. Moler와 Van Loan은 1978년의 논문에서 행렬 지수를 계산하는 열아홉 가지 방법을 비교하며 어느 것에나 문제가 있음을 보였고 25년 뒤의 재검토에서도 결론이 크게 바뀌지 않았다\cite{moler2003nineteen}. 오늘날에는 $e^A=(e^{A/2^s})^{2^s}$로 노름을 줄인 뒤 Padé 근사를 쓰고 제곱을 되풀이하는 방법이 표준이다. stardust의 \code{linalg.expm}이 이 방법을 쓰며 57장에서 오차 분석과 함께 다룬다.

\cref{fig:airy}에 이계 방정식의 두 해가 얼마나 다르게 움직일 수 있는지 그렸다. Airy 방정식 $y''=xy$는 $x<0$에서 $y''=-|x|y$이므로 진동하고 $x>0$에서는 지수적으로 자라거나 준다. 두 표준 해 가운데 $\operatorname{Ai}$는 $x\to\infty$에서 줄고 $\operatorname{Bi}$는 자란다. $p=0$이므로 Wronskian은 상수이고 그 값은 $W(\operatorname{Ai},\operatorname{Bi})=1/\pi$이다. 따로 빌드한 stardust의 \code{special.airy}가 낸 네 값으로 $\pi W-1$을 재면 $[-10,2.5]$의 $501$점에서 최대 $1.8\times10^{-15}$이었다. 점화식에서도 같은 두 행동, 곧 진동하는 영역과 한 해가 다른 해를 압도하는 영역이 나타난다.

\begin{figure}[htbp]
\centering
\begin{tikzpicture}
\begin{axis}[sdplot,xlabel={$x$},ylabel={$y$},xmin=-10,xmax=2.5,ymin=-0.6,ymax=1.5]
\addplot[black,thick] table[x=x,y=ai]{figures/06/airy.dat};
\addlegendentry{$\operatorname{Ai}(x)$}
\addplot[black,dashed,thick] table[x=x,y=bi]{figures/06/airy.dat};
\addlegendentry{$\operatorname{Bi}(x)$}
\end{axis}
\end{tikzpicture}
\caption{Airy 방정식 $y''=xy$의 두 해(따로 빌드한 stardust \code{special.airy}; \code{figures/06/longrec.py}). 전이점 $x=0$의 왼쪽은 진동, 오른쪽은 지수적 분리.}
\label{fig:airy}
\end{figure}


\section{Equations of Special Functions}
\label{sec:special-odes}

\subsection{The Classical Equations}

계수가 다항식이나 유리 함수인 이계 방정식의 해는 대개 초등 함수로 나타나지 않는다. 물리학과 공학에서 되풀이해 나오는 몇몇 방정식의 해에는 이름이 붙었다.
\begin{itemize}
\item \textbf{Bessel 방정식} $x^2y''+xy'+(x^2-\nu^2)y=0$. 원기둥 좌표에서 파동방정식이나 열방정식을 변수분리하면 나온다. 해는 Bessel 함수 $J_\nu$, $Y_\nu$이다. \cref{thm:abel}에서 $p=1/x$이므로 Wronskian은 $C/x$ 꼴이고 실제로 $W(J_\nu,Y_\nu)=2/(\pi x)$이다\cite{watson1944bessel}.
\item \textbf{변형 Bessel 방정식} $x^2y''+xy'-(x^2+\nu^2)y=0$. 해는 지수적으로 자라는 $I_\nu$와 줄어드는 $K_\nu$이고 $W(K_\nu,I_\nu)=1/x$이다.
\item \textbf{Airy 방정식} $y''=xy$. 빛의 무지개 근처 세기와 양자역학의 선형 퍼텐셜처럼 전이 영역에서 나온다.
\item \textbf{초기하 방정식} $x(1-x)y''+(c-(a+b+1)x)y'-aby=0$. 정칙 특이점이 $0,1,\infty$ 셋뿐인 이계 방정식은 모두 이 꼴로 바뀌므로 많은 특수 함수가 그 특수한 경우이다.
\end{itemize}
\cref{tab:special-odes}에 이 방정식들의 Wronskian과 차수에 대한 점화식을 모았다. Wronskian의 상수는 $x\to0$이나 $x\to\infty$에서 두 해의 주항을 넣어 구한다.

\begin{table}[htbp]
\begin{sdtable}{colspec={llll}, row{1}={font=\small\bfseries}}
\toprule
해 & Wronskian & 차수 점화식 & 최소해 / 지배해 ($n\to\infty$) \\
\midrule
$J_\nu$, $Y_\nu$ & $W(J_\nu,Y_\nu)=2/(\pi x)$ & $y_{\nu+1}=(2\nu/x)y_\nu-y_{\nu-1}$ & $J$ / $Y$ \\
$I_\nu$, $K_\nu$ & $W(K_\nu,I_\nu)=1/x$ & $I_{\nu+1}=-(2\nu/x)I_\nu+I_{\nu-1}$, $K_{\nu+1}=(2\nu/x)K_\nu+K_{\nu-1}$ & $I$ / $K$ \\
$\operatorname{Ai}$, $\operatorname{Bi}$ & $W(\operatorname{Ai},\operatorname{Bi})=1/\pi$ & (차수 없음) & -- \\
\bottomrule
\end{sdtable}
\caption{Bessel·Airy 방정식의 해의 Wronskian과 점화식\cite{watson1944bessel,dlmf}.}
\label{tab:special-odes}
\end{table} 예컨대 $x\to\infty$에서 $J_\nu\sim\sqrt{2/(\pi x)}\cos\chi$, $Y_\nu\sim\sqrt{2/(\pi x)}\sin\chi$($\chi=x-\nu\pi/2-\pi/4$)이고 $W$는 $(2/(\pi x))(\cos^2\chi+\sin^2\chi)=2/(\pi x)$이다.

\subsection{Series Solutions}

이런 방정식의 해를 원점 근처에서 구하는 표준적인 방법은 $y=\sum_kc_kx^{k+r}$ 꼴의 급수를 넣어 계수의 점화식을 얻는 Frobenius 방법이다. Bessel 방정식에 넣으면 $x^{k+r}$의 계수를 비교해 $((k+r)^2-\nu^2)c_k+c_{k-2}=0$을 얻고 가장 낮은 차수에서 $r=\pm\nu$가 정해진다. $r=\nu$로 두고 정규화하면
\begin{equation}
\label{eq:jv-series}
J_\nu(x)=\Bigl(\frac x2\Bigr)^\nu\sum_{k=0}^\infty\frac{(-x^2/4)^k}{k!\,\Gamma(k+\nu+1)}
\end{equation}
을 얻는다. 이 급수는 모든 $x$에서 수렴하지만 $x$가 크면 항들이 크게 자랐다가 서로 지워진다. 가장 큰 항은 $(x/2)^{2k}/(k!)^2$이 최대인 $k\approx x/2$에서 대략 $e^x/(\pi x)$이다. $x=30$이면 약 $10^{11}$인 항들이 더해져 $|J_0(30)|\approx0.09$가 되므로 $10^{12}$배의 상쇄가 생긴다. 수치 계산에는 작은 $x$에서만 쓴다. stardust의 \code{jv}는 $x^2<2(\nu+1)$에서만 이 급수를 쓴다(\code{csrc/special/ops/bessel.hpp}의 \code{fast\_j}). 그 범위에서는 연속한 항의 비 $x^2/(4k(\nu+k))$가 $1/2$ 아래이므로 항이 한 번에 절반 이하로 줄고 교대합의 상쇄는 세 배를 넘지 않는다. $\Gamma(\nu+1)$은 $\nu\ge1$이면 $\nu\Gamma(\nu)$로 계산한다. $\nu+1$을 그대로 쓰면 그 값이 반올림될 때(이진 지수의 경계를 넘을 때) 상대오차가 약 $\psi(\nu+1)(\nu+1)/2$ ulp로 커진다. 큰 $x$에서는 Hankel 점근 전개(92장; 3장의 \cref{alg:adaptive-trunc}처럼 가장 작은 항에서 자른다)를, 그 사이에서는 다음 절부터의 점화식과 연분수를 쓴다.

Frobenius 방법은 다음 조건에서 통한다.

\begin{theorem}[Frobenius]
\label{thm:frobenius}
$x^2y''+xP(x)y'+Q(x)y=0$에서 $P$, $Q$가 $|x|<R$에서 수렴하는 거듭제곱 급수이고 지표 방정식 $F(r)=r(r-1)+P(0)r+Q(0)=0$의 두 근이 실수 $r_1\ge r_2$라 하자. 그러면 $y_1=x^{r_1}\sum_kc_kx^k$($c_0=1$)은 $0<x<R$에서 수렴하는 해이다. $r_1-r_2$가 정수가 아니면 $x^{r_2}\sum_kd_kx^k$가 둘째 해이고 정수이면 둘째 해는 $Cy_1\ln x+x^{r_2}\sum_kd_kx^k$ 꼴이다.
\end{theorem}

\begin{proof}
급수를 넣고 $x^{r+k}$의 계수를 비교하면 $F(r+k)c_k=-\sum_{j<k}((r+j)P_{k-j}+Q_{k-j})c_j$이다. $r=r_1$이면 $F(r_1+k)=k(k+r_1-r_2)\ne0$($k\ge1$)이므로 계수가 하나씩 정해진다. $\rho<R$에서 $|P_j|,|Q_j|\le M\rho^{-j}$이면 $|F(r_1+k)|\ge k^2$를 써서 귀납적으로 $|c_k|\le C\rho'^{-k}$($\rho'<\rho$는 아무것이나 좋다)를 보일 수 있으므로 급수가 수렴한다. 5장 \cref{cor:cauchy-estimates}의 Cauchy 부등식과 같은 비교이다. 수렴하는 급수를 미분방정식에 넣으면 계수 관계가 곧 방정식이므로 해이다. $r_2$에서는 $F(r_2+k)=0$이 $k=r_1-r_2$에서만 일어난다. 그때 계수 방정식을 풀 수 없으면 $c_k(r)$을 $r$의 함수로 보고 $\partial/\partial r$을 취해 로그 항을 만든다. 자세한 것은 Coddington과 Levinson\cite{coddington1955}(4장)에 있다.
\end{proof}

Bessel 방정식은 $P=1$, $Q=x^2-\nu^2$이므로 지표 근이 $\pm\nu$이다. $\nu$가 정수가 아니면 $J_\nu$와 $J_{-\nu}$가 기저이다($2\nu$가 홀수인 반정수에서는 \cref{thm:frobenius}의 로그 항이 나타나지 않는다). 둘째 종류의 함수는
\begin{equation}
\label{eq:yv-def}
Y_\nu(x)=\frac{J_\nu(x)\cos\nu\pi-J_{-\nu}(x)}{\sin\nu\pi}
\end{equation}
로 정의한다. 정수 차수에서는 분자와 분모가 함께 $0$이 되므로 극한을 취하고 그 극한에서 \cref{thm:frobenius}의 로그 항이 나온다. 이 정의를 그대로 계산하면 $\nu$가 정수 근처일 때 상쇄가 생긴다. Temme\cite{temme1975bessel}는 $|\mu|\le1/2$에서 이 극한을 미리 풀어 둔 급수를 내놓았다.

\begin{proposition}[Temme의 급수]
\label{prop:temme-series}
$|\mu|\le1/2$, $x>0$에서 $c_k=(-x^2/4)^k/k!$,
\[
\Gamma_1(\mu)=\frac{1}{2\mu}\Bigl(\frac{1}{\Gamma(1-\mu)}-\frac{1}{\Gamma(1+\mu)}\Bigr),\qquad\Gamma_2(\mu)=\frac12\Bigl(\frac{1}{\Gamma(1-\mu)}+\frac{1}{\Gamma(1+\mu)}\Bigr),
\]
$\sigma=\mu\ln(2/x)$, $p_0=\pi^{-1}(x/2)^{-\mu}\Gamma(1+\mu)$, $q_0=\pi^{-1}(x/2)^\mu\Gamma(1-\mu)$,
\[
f_0=\frac2\pi\cdot\frac{\mu\pi}{\sin\mu\pi}\Bigl(\Gamma_1(\mu)\cosh\sigma+\Gamma_2(\mu)\ln\frac2x\cdot\frac{\sinh\sigma}{\sigma}\Bigr)
\]
로 두고 $p_k=p_{k-1}/(k-\mu)$, $q_k=q_{k-1}/(k+\mu)$, $f_k=(kf_{k-1}+p_{k-1}+q_{k-1})/(k^2-\mu^2)$라 하고 $g_k=f_k+(2/\mu)\sin^2(\mu\pi/2)\,q_k$로 두면
\[
Y_\mu(x)=-\sum_{k\ge0}c_kg_k,\qquad Y_{\mu+1}(x)=-\frac2x\sum_{k\ge0}c_k(p_k-kg_k)
\]
이다. $\mu\to0$에서 각 계수는 연속으로 이어진다.
\end{proposition}

증명은 줄거리만 적는다. \cref{eq:yv-def}에 \cref{eq:jv-series}와 $J_{-\mu}$의 급수를 넣고 $c_k$마다 묶으면 $k$번째 계수는 $(x/2)^{\pm\mu}/\Gamma(k+1\pm\mu)$들의 조합을 $\sin\mu\pi$로 나눈 꼴이 된다. 이것을 $\cosh\sigma$, $\sinh\sigma$와 $1/\Gamma(1\pm\mu)$의 합과 차로 다시 쓰면 $\mu\to0$에서 $0/0$인 부분이 모두 $\Gamma_1$, $\Gamma_2$, $\sinh\sigma/\sigma$, $\mu\pi/\sin\mu\pi$처럼 $\mu=0$에서 매끄러운 함수로 모인다. $f_k$, $p_k$, $q_k$의 점화식은 $1/\Gamma(k+1\pm\mu)=1/((k\pm\mu)\Gamma(k\pm\mu))$에서 나온다. 전 과정은 Temme의 원 논문과 Numerical Recipes\cite{press2007nr}(§6.6)에 있다. 항이 $c_k\sim(x^2/4)^k/k!$처럼 줄므로 이 급수는 $x$가 $2$ 근처까지 빠르고 정확하다. $x$가 크면 $c_k$가 자랐다가 줄며 상쇄가 생긴다. stardust는 $\Gamma_1$, $\Gamma_2$를 Numerical Recipes의 맞춘 Chebyshev 급수 대신 $1/\Gamma(1+z)$의 Taylor 급수로 계산한다.

\subsection{Turning Points}

$y''=q(x)y$ 꼴의 방정식에서 $q>0$인 곳은 해가 지수적으로 자라거나 줄고 $q<0$인 곳은 진동한다. $q$가 천천히 변하면 근사해 $q^{-1/4}\exp(\pm\int\sqrt q)$(Liouville--Green 또는 WKB 근사)로 이 구분을 정량적으로 쓸 수 있다. 그 오차 한계와 $q$의 부호가 바뀌는 \term{전이점}{turning point} 근처에서 Airy 함수로 이어 붙이는 방법은 Olver\cite{olver1974asymptotics}(6장, 11장)가 정리했다. 증명은 대략 이렇다. 변수 바꾸기로 방정식을 $w''=(1+\text{작은 항})w$ 꼴로 바꾸고 작은 항을 Volterra 적분 방정식으로 다루어 \cref{thm:ode-exist}의 Picard 반복처럼 한계를 얻는다. Bessel 방정식을 $y=x^{-1/2}w$로 바꾸면 $w''=((\nu^2-1/4)/x^2-1)w$이고 전이점은 $x\approx\nu$이다. $J_\nu(x)$가 $x>\nu$에서 진동하고 $x<\nu$에서 지수적으로 작은 것도 이 때문이다. 그 크기는 Liouville--Green 근사의 지수가 정한다. $\nu^2\gg1/4$이면 $q(t)\approx\nu^2/t^2-1$이다.

\begin{proposition}
\label{prop:debye-exponent}
$0<x<\nu$, $\operatorname{sech}\alpha=x/\nu$이면 $\int_x^\nu\sqrt{\nu^2/t^2-1}\,dt=\nu(\alpha-\tanh\alpha)$이다.
\end{proposition}

\begin{proof}
$t=\nu\operatorname{sech}\theta$로 바꾸면 $\sqrt{\nu^2/t^2-1}=\sqrt{\cosh^2\theta-1}=\sinh\theta$, $dt=-\nu\operatorname{sech}\theta\tanh\theta\,d\theta$이고 $t=x$는 $\theta=\alpha$, $t=\nu$는 $\theta=0$이다. 그러므로 적분은 $\nu\int_0^\alpha\tanh^2\theta\,d\theta=\nu\int_0^\alpha(1-\operatorname{sech}^2\theta)\,d\theta=\nu(\alpha-\tanh\alpha)$이다.
\end{proof}

전이점 $t=\nu$에서 크기가 $1$ 정도인 감쇠 해는 $x$에서 $e^{-\nu(\alpha-\tanh\alpha)}$ 정도로 작고 자라는 해는 그 역수만큼 크다. 이것이 4절에서 쓰는 Debye 점근식 $J_\nu(\nu\operatorname{sech}\alpha)\approx e^{-\nu(\alpha-\tanh\alpha)}/\sqrt{2\pi\nu\tanh\alpha}$의 지수이다. 앞의 인자는 Liouville--Green 근사의 $q^{-1/4}$와 전이점에서 Airy 함수로 이어 붙인 상수에서 나온다(\cite{dlmf} 10.19.3). stardust는 이 지수를 두 군데에서 쓴다. 하나는 다음 절의 시작 차수이고 다른 하나는 $J_\nu(x)$가 아래로 넘쳐 $0$이 되는지 미리 판정하는 \code{j\_underflows}이다. 판정에는 $\nu\ge-1/2$에서 성립하는 $|J_\nu(x)|\le(x/2)^\nu/\Gamma(\nu+1)$ 같은 엄밀한 상한(\cite{dlmf} §10.14)을 쓴다. 같은 구조가 차수에 대한 점화식에서 $n<x$와 $n>x$의 구분으로 다시 나타난다(\cref{fig:bessel-mag}).


\section{Three-Term Recurrences}
\label{sec:three-term}

\subsection{The Solution Space}

Bessel 함수는 차수에 대해
\begin{equation}
\label{eq:bessel-rec}
y_{n+1}=\frac{2n}{x}y_n-y_{n-1}
\end{equation}
을 만족한다. $J_n(x)$와 $Y_n(x)$가 모두 이 점화식의 해이고 변형 Bessel 함수는 $I_{n+1}=-(2n/x)I_n+I_{n-1}$, $K_{n+1}=(2n/x)K_n+K_{n-1}$을 만족한다. 이처럼 이웃한 세 항을 잇는 선형 점화식
\begin{equation}
\label{eq:three-term}
y_{n+1}+a_ny_n+b_ny_{n-1}=0\qquad(b_n\ne0)
\end{equation}
을 \term{삼항 점화식}{three-term recurrence}이라 한다. 처음 두 값이 나머지를 모두 정하고 $b_n\ne0$이면 거꾸로도 정하므로 해 공간은 $2$차원이다. Wronskian의 이산판은 \term{Casorati 행렬식}{Casoratian} $C_n=y_nz_{n+1}-y_{n+1}z_n$이다.

\begin{proposition}
\label{prop:casorati}
\cref{eq:three-term}의 두 해에 대해 $C_n=b_nC_{n-1}$이다. 따라서 $C_n=C_0\prod_{k=1}^nb_k$이고 두 해가 독립일 필요충분조건은 $C_0\ne0$이다.
\end{proposition}

\begin{proof}
$y_{n+1}=-a_ny_n-b_ny_{n-1}$과 $z_{n+1}$의 같은 식을 넣으면 $C_n=y_n(-a_nz_n-b_nz_{n-1})-(-a_ny_n-b_ny_{n-1})z_n=b_n(y_{n-1}z_n-y_nz_{n-1})=b_nC_{n-1}$이다. $b_n\ne0$이므로 $C_n$은 모두 $0$이거나 모두 $0$이 아니다. $C_0=0$이면 초깃값 $(y_0,y_1)$과 $(z_0,z_1)$이 비례하고 해가 초깃값으로 정해지므로 두 해가 비례한다.
\end{proof}

Bessel 점화식 \cref{eq:bessel-rec}은 $b_n=1$이므로 $C_n$이 상수이다. 실제로 Wronskian 관계를 점화식으로 바꾸면 $J_{n+1}Y_n-J_nY_{n+1}=2/(\pi x)$이다.

\subsection{Constant Coefficients}

\begin{proposition}
\label{prop:const-coeff}
상수 $a,b$($b\ne0$)에 대해 $y_{n+1}+ay_n+by_{n-1}=0$의 특성방정식 $\lambda^2+a\lambda+b=0$의 두 근이 $\lambda_1\ne\lambda_2$이면 일반해는 $y_n=c_1\lambda_1^n+c_2\lambda_2^n$이다. 중근 $\lambda$이면 $y_n=(c_1+c_2n)\lambda^n$이다.
\end{proposition}

\begin{proof}
$y_n=\lambda^n$을 넣으면 $\lambda^{n-1}(\lambda^2+a\lambda+b)=0$이므로 특성근의 거듭제곱은 해이다. 두 근이 다르면 \cref{prop:casorati}의 $C_0=\lambda_2-\lambda_1\ne0$이므로 두 해가 기저이다. 중근이면 $a=-2\lambda$, $b=\lambda^2$이고 $n\lambda^n$을 넣으면 $\lambda^{n-1}((n+1)\lambda^2-2n\lambda^2+(n-1)\lambda^2)=0$이다. $C_0=\lambda\ne0$이다.
\end{proof}

특성근은 복소평면의 점이고 해의 크기는 근의 절댓값으로 정해진다. 두 근의 절댓값이 같으면($|\lambda_1|=|\lambda_2|$, 예컨대 켤레 복소근) 두 해가 같은 속도로 자라며 진동한다. 다르면 큰 근의 거듭제곱이 작은 근의 거듭제곱을 압도한다. 이 둘이 Airy 방정식의 진동 영역과 지수 영역에 해당한다. 계수가 $n$에 따라 변하면 특성근도 $n$마다 달라진다. Bessel 점화식의 특성방정식 $\lambda^2-(2n/x)\lambda+1=0$은 $n<x$에서 절댓값 $1$인 켤레 복소근을, $n>x$에서 곱이 $1$인 두 실근 $\lambda_\pm=n/x\pm\sqrt{(n/x)^2-1}$을 갖는다. $J_n$과 $Y_n$은 $n<x$에서 비슷한 크기로 진동하다가 $n>x$에서 갈라진다(\cref{fig:bessel-mag}). 계수가 수렴하면 이 관찰을 정리로 세울 수 있다. Poincaré--Perron 정리에 따르면 $a_n\to a$, $b_n\to b$이고 극한 특성방정식의 두 근의 절댓값이 다르면 자명하지 않은 모든 해의 비 $y_{n+1}/y_n$이 충분히 큰 $n$에서 정의되고 어느 한 근으로 수렴한다. 증명에서는 비 $r_n=y_{n+1}/y_n$이 만족하는 Riccati 꼴 점화식 $r_n=-a_n-b_n/r_{n-1}$을 두 근 근처의 축소 사상으로 본다. 자세한 내용은 Gautschi의 해설\cite{gautschi1967recurrence}(§1)과 그 참고문헌에 있다. Bessel 점화식은 $a_n=-2n/x$가 발산하므로 이 정리가 그대로 적용되지 않는다. Perron과 Kreuser는 이 정리를 계수가 $n$의 거듭제곱처럼 자라는 경우로 넓혔지만 Bessel 함수에서는 그럴 필요 없이 $J_n$과 $Y_n$의 점근식(아래)으로 직접 판정한다.

\begin{figure}[htbp]
\centering
\begin{tikzpicture}
\begin{semilogyaxis}[sdplot,xlabel={차수 $n$},ylabel={크기},xmin=0,xmax=40,ymin=1e-16,ymax=1e25,ytick={1e-15,1e-10,1e-5,1,1e5,1e10,1e15,1e20,1e25}]
\addplot[black,mark=*,mark size=1.1pt] table[x=n,y=J]{figures/06/bessel_mag.dat};
\addlegendentry{$|J_n(10)|$}
\addplot[black,dashed,mark=square*,mark size=1.1pt] table[x=n,y=Y]{figures/06/bessel_mag.dat};
\addlegendentry{$|Y_n(10)|$}
\addplot[black!40,thin,forget plot] coordinates {(10,1e-16) (10,1e25)};
\end{semilogyaxis}
\end{tikzpicture}
\caption{$x=10$에서 차수에 따른 $|J_n|$과 $|Y_n|$(따로 빌드한 stardust; \code{figures/06/bessel.py}). $n<x$에서는 둘 다 $\sqrt{2/(\pi x)}$ 근처에서 진동하고 $n>x$에서는 $J_n$이 줄고 $Y_n$이 자란다.}
\label{fig:bessel-mag}
\end{figure}

\subsection{Minimal and Dominant Solutions}

\begin{definition}[최소해와 지배해]
\label{def:minimal}
\cref{eq:three-term}의 해 $f_n$이 그것과 독립인 모든 해 $g_n$에 대해 $f_n/g_n\to0$($n\to\infty$)을 만족하면 $f$를 \term{최소해}{minimal solution}라 하고 그런 $g$를 \term{지배해}{dominant solution}라 한다.
\end{definition}

최소해가 있다면 상수배를 빼고 유일하다. 두 최소해 $f$, $f'$가 독립이면 정의에서 $f/f'\to0$과 $f'/f\to0$이 함께 성립해야 하므로 모순이다. \cref{prop:const-coeff}에서 $|\lambda_1|<|\lambda_2|$이면 $\lambda_1^n$이 최소해이다. $c_2\ne0$이면 $\lambda_1^n/(c_1\lambda_1^n+c_2\lambda_2^n)\to0$이기 때문이다. Bessel 점화식에서는 고정된 $x$와 $n\to\infty$에 대해
\[
J_n(x)\sim\frac{(x/2)^n}{n!},\qquad Y_n(x)\sim-\frac{(n-1)!}{\pi}\Bigl(\frac2x\Bigr)^n
\]
이므로(\cref{eq:jv-series}의 첫 항과 $Y_n$의 급수의 첫 항) $J_n$이 최소해이고 $Y_n$이 지배해이다. 변형 Bessel 점화식에서는 $I_n$이 최소해, $K_n$이 지배해이다.

\subsection{Direction and Stability}

최소해를 앞으로 계산하면 무슨 일이 생기는지 보자. 계산 중에 생긴 오차가 퍼지는 모양은 미분방정식의 상수 변화법을 이산판으로 옮긴 다음 명제로 정확히 쓸 수 있다.

\begin{proposition}[이산 상수 변화법]
\label{prop:discrete-green}
$f$, $g$가 \cref{eq:three-term}의 독립인 두 해이고 $C_k=f_kg_{k+1}-f_{k+1}g_k$라 하자. 앞으로 계산하는 동안 $k$번째 걸음($y_{k+1}$을 만드는 걸음)에 오차 $\delta_k$가 더해지면 $n>k$에서의 오차는
\[
e_n=\sum_{k<n}\delta_k\frac{f_kg_n-g_kf_n}{C_k}
\]
이다. 거꾸로 계산하며 $y_{k-1}$을 만드는 걸음의 오차 $\delta'_k$는 $n<k$에서 $\sum_{k>n}\delta'_k(f_kg_n-g_kf_n)/(f_kg_{k-1}-g_kf_{k-1})$로 퍼진다.
\end{proposition}

\begin{proof}
점화식이 선형이므로 각 걸음의 오차는 따로 퍼져 더해진다. 걸음 $k$에서 들어온 오차 하나만 보면 그 오차는 $e_k=0$, $e_{k+1}=\delta_k$에서 출발하는 동차 점화식의 해이다. $h_n=f_kg_n-g_kf_n$은 해이고 $h_k=0$, $h_{k+1}=C_k$이므로 그 해는 $\delta_kh_n/C_k$이다. 거꾸로 가는 경우도 같다.
\end{proof}

이 명제에서 방향의 규칙이 모두 나온다. 최소해 $f$를 앞으로 계산하면 걸음마다의 오차 $\delta_k\approx u|f_{k+1}|$가 $n$에서
\[
\frac{e_n}{f_n}\approx\sum_k\delta_k\frac{f_k}{C_k}\cdot\frac{g_n}{f_n}
\]
를 낸다($g_kf_n$ 항은 상대적으로 작다). 비 $g_n/f_n$이 $\infty$로 가므로 상대오차가 함께 자란다. 지배해 $g$를 앞으로 계산하면 $\delta_k\approx u|g_{k+1}|$이고 $e_n/g_n\approx\sum_k\delta_kf_k/C_k$이다. 최소해와 지배해가 갈라진 영역에서는 $C_k\approx f_kg_{k+1}$이므로 각 항이 $u$ 정도이다. 오차는 증폭되지 않고 걸음마다 $u$씩 더해져 $n$걸음 뒤에 $nu$(같은 부호로 쌓일 때)나 $\sqrt nu$(무작위일 때) 정도가 된다.

\begin{corollary}
\label{prop:forward-error}
최소해를 앞으로 계산한 값의 상대오차는 $g_n/f_n$에 비례해 자란다. 지배해를 앞으로 계산하거나 최소해를 거꾸로 계산하면 상대오차는 걸음 수에 비례하는 정도로만 쌓인다.
\end{corollary}

최소해는 앞으로 계산하면 불안정하지만 지배해는 앞으로 계산해도 안정하다. 지배해에 섞여 드는 최소해 성분은 상대적으로 작아지기 때문이다. 예를 \cref{fig:bessel-dir}에 그렸다. 십진 산술 $90$자리로 구한 $J_0(1)$과 $J_1(1)$에서 출발해 \cref{eq:bessel-rec}를 binary64로 앞으로 돌리면 상대오차가 $n=5$에서 $6.2\times10^{-11}$, $n=10$에서 $27$이다. $n=20$에서는 참값 $3.9\times10^{-25}$ 대신 $-2.4\times10^6$이 나온다. 점화식은 맞는데 돌린 방향이 틀려서 계산이 망가졌다. 이 현상과 처방은 Gautschi의 해설\cite{gautschi1967recurrence}이 체계적으로 정리했다. 같은 그림의 Miller 방법과 stardust의 \code{jv}는 모든 $n$에서 $3.1\times10^{-16}$ 아래이다.

\begin{figure}[htbp]
\centering
\begin{tikzpicture}
\begin{semilogyaxis}[sdplot,xlabel={차수 $n$},ylabel={$J_n(1)$의 상대오차},xmin=0,xmax=25,ymin=1e-17,ymax=1e4,ytick={1e-16,1e-12,1e-8,1e-4,1},yticklabels={${\le}10^{-16}$,$10^{-12}$,$10^{-8}$,$10^{-4}$,$1$}]
\addplot[black,mark=*,mark size=1.2pt] table[x=n,y=fwd]{figures/06/bessel_dir.dat};
\addlegendentry{앞으로 점화}
\addplot[black,dashed,mark=square*,mark size=1.2pt] table[x=n,y=miller]{figures/06/bessel_dir.dat};
\addlegendentry{Miller ($N=30$)}
\addplot[black!50,mark=triangle*,mark size=1.4pt] table[x=n,y=lib]{figures/06/bessel_dir.dat};
\addlegendentry{stardust \code{jv}}
\end{semilogyaxis}
\end{tikzpicture}
\caption{$J_n(1)$을 세 방법으로 계산한 상대오차(\code{figures/06/bessel.py}). 앞으로 점화의 값 $10^3$은 그 이상을 뜻한다.}
\label{fig:bessel-dir}
\end{figure}

$n<x$인 영역에서는 두 해가 같은 크기로 진동하므로 어느 방향도 다른 쪽을 압도하지 않는다. 이 영역에서 점화는 \term{중립}{neutral}이어서 반올림이 쌓이기만 하고 증폭되지는 않는다. stardust의 \code{jv}는 $\nu<x$이고 $J_\mu$, $J_{\mu+1}$($\mu=\nu-\operatorname{round}(\nu)$)을 싸게 얻을 수 있으면(차수 $0$과 $1$의 유리 근사, $\mu=-1/2$의 닫힌 꼴, $x\ge25$의 Hankel 전개) 위로 점화한다. 그 밖에서는 다음 절의 역방향 점화를 쓴다.


\section{Backward Recurrence and Continued Fractions}
\label{sec:miller}

\subsection{Miller's Algorithm}

최소해를 안정하게 계산하려면 점화식을 \textbf{거꾸로} 돌리면 된다. 거꾸로 가면 역할이 바뀌어 최소해가 지배적으로 자란다. 그런데 거꾸로 시작하려면 큰 $N$의 두 값이 필요한데 바로 그 값을 모른다. Miller의 해법은 단순하다. 충분히 큰 $N$에서 아무 값, 예컨대 $\tilde f_{N+1}=0$과 $\tilde f_N=1$로 시작해 거꾸로 계산하고 마지막에 알려진 관계식으로 전체의 척도를 맞춘다(\cref{alg:miller}).

\begin{algorithm}[htbp]
\begin{algorithmic}[1]
\Require 점화식 \cref{eq:three-term}, 출발 차수 $N$, 정규화 관계 $\sum_k\mu_kf_k=s$
\Ensure 최소해의 근삿값 $f_0,\dots,f_n$ ($n<N$)
\State $\tilde f_{N+1}\gets0$, $\tilde f_N\gets1$
\For{$k\gets N$ \textbf{downto} $1$}
  \State $\tilde f_{k-1}\gets-(\tilde f_{k+1}+a_k\tilde f_k)/b_k$ \Comment{점화식을 $y_{k-1}$에 대해 푼 것}
  \State \textbf{if} $|\tilde f_{k-1}|>2^{600}$ \textbf{then} 지금까지의 $\tilde f$와 부분합에 $2^{-600}$을 곱한다
\EndFor
\State $\sigma\gets\sum_k\mu_k\tilde f_k$ \Comment{거꾸로 가며 함께 누적할 수 있다}
\For{$k\gets0$ \textbf{to} $n$} $f_k\gets s\,\tilde f_k/\sigma$
\EndFor
\end{algorithmic}
\caption{Miller의 역방향 점화}
\label{alg:miller}
\end{algorithm}

\begin{theorem}
\label{thm:miller}
\cref{eq:three-term}이 최소해 $f$와 지배해 $g$를 갖고 $f_0\ne0$이라 하자. \cref{alg:miller}의 정규화 전 수열을 $\tilde f^{(N)}$이라 하면 고정된 $k$에 대해 $N\to\infty$에서 $\tilde f^{(N)}_k/\tilde f^{(N)}_0\to f_k/f_0$이고 그 오차는 첫 차수에서
\[
\frac{\tilde f^{(N)}_k/\tilde f^{(N)}_0}{f_k/f_0}-1\approx-\frac{f_{N+1}}{g_{N+1}}\Bigl(\frac{g_k}{f_k}-\frac{g_0}{f_0}\Bigr)
\]
이다.
\end{theorem}

\begin{proof}
$\tilde f^{(N)}$은 \cref{eq:three-term}의 해이므로 $\tilde f^{(N)}=\alpha_Nf+\beta_Ng$로 쓸 수 있다. 조건 $\tilde f^{(N)}_{N+1}=0$에서 $\beta_N/\alpha_N=-f_{N+1}/g_{N+1}=:-\epsilon_N$이다. 따라서
\[
\frac{\tilde f^{(N)}_k}{\tilde f^{(N)}_0}=\frac{f_k-\epsilon_Ng_k}{f_0-\epsilon_Ng_0}=\frac{f_k}{f_0}\cdot\frac{1-\epsilon_Ng_k/f_k}{1-\epsilon_Ng_0/f_0}
\]
이다. 최소해의 정의에서 $\epsilon_N\to0$이므로 $k$를 고정하면 비는 $f_k/f_0$으로 가고 $\epsilon_N$에 대해 전개하면 오차의 식이 나온다.
\end{proof}

정리에 따라 거꾸로 계산한 수열은 \textbf{비율}이 정확하므로 척도 하나만 맞추면 된다. 오차의 크기는 $\epsilon_N=f_{N+1}/g_{N+1}$과 $g_k/f_k$의 곱이다. 앞쪽 인자는 $N$이 클수록 작아지고 뒤쪽 인자는 원하는 차수 $k$가 클수록 커지므로 $N$은 원하는 가장 큰 차수보다 충분히 커야 한다. 정확한 산술로도 생기는 이 오차와 별개로 반올림도 따져야 한다. 거꾸로 가는 동안 각 단계의 반올림은 그 자리에서 시작하는 해를 더하는 것과 같다. 거꾸로 가는 방향에서는 최소해가 지배적이므로 그 오차들은 증폭되지 않고 걸음마다 $u$ 정도씩 더해질 뿐이다(\cref{sec:dominant}의 긴 점화 참고). 이 방법은 J. C. P. Miller가 1950년대 영국 학술협회의 Bessel 함수 수표를 만들며 쓴 데서 이름이 붙었다\cite{gautschi1967recurrence}.

\subsection{Choosing the Starting Order}

Bessel 함수에서는 $\epsilon_N=J_{N+1}/Y_{N+1}$을 차수 $N>x$에서 얼마나 작게 할지 정해야 한다. Debye의 점근 전개(\cite{dlmf} 10.19.3)에 따르면 $N>x$에서 $\operatorname{sech}\alpha=x/N$으로 두면
\[
J_N(x)\approx\frac{e^{-N(\alpha-\tanh\alpha)}}{\sqrt{2\pi N\tanh\alpha}},\qquad Y_N(x)\approx-\frac{e^{N(\alpha-\tanh\alpha)}}{\sqrt{\pi N\tanh\alpha/2}}
\]
이다. 두 식을 나누면 $|\epsilon_N|\approx e^{-2N(\alpha-\tanh\alpha)}/2$이고 이것이 $2^{-53}$ 아래로 내려가려면 $N(\alpha-\tanh\alpha)$가 약 $18$을 넘어야 한다. 여유를 두어 $19$로 잡자. 원하는 차수가 전이점 $x$ 근처이면 $g_k/f_k$가 $1$ 근처이므로 이 여유로 충분하다. 전이점 근처에서 $\alpha$가 작으면 $\alpha-\tanh\alpha\approx\alpha^3/3$이고 $\alpha^2\approx2(N-x)/N$이다. 따라서 조건은
\[
\frac N3\Bigl(\frac{2(N-x)}{N}\Bigr)^{3/2}\gtrsim19,\qquad\text{곧}\qquad N-x\gtrsim7.4\,N^{1/3}
\]
이다. 필요한 여유는 차수의 세제곱근으로만 자란다. stardust의 \code{miller}(\code{csrc/special/ops/bessel.hpp})는 $\mathrm{top}=\max(\nu,x)$에 대해 $N\ge\mathrm{top}+10\,\mathrm{top}^{1/3}+8$을 쓴다. 코드의 주석에 따르면 $\mathrm{top}$이 $0.5$--$100$일 때 $7\,\mathrm{top}^{1/3}+3$이면 비 $f_0/f_1$이 자리를 잡는다. 그런데 정규화에 쓰는 Neumann 합(아래)이 $\mathrm{top}+2\,\mathrm{top}^{1/3}$까지의 차수를 쓰고 그 오차 $\epsilon_NY_k$가 $k$와 함께 자라므로 여유를 더 두었다. \cref{fig:miller-n}에서 이 여유를 다시 쟀다. $\mu=0.3$, $x=10$에서 정규화한 $J_{0.3}$과 $J_{20.3}$의 오차가 $N-\mathrm{top}$에 따라 지수적으로 줄어 각각 $24$와 $14$ 근처에서 반올림 수준에 닿는다. stardust의 하한 $10\,\mathrm{top}^{1/3}+8$은 $29.5$와 $35.3$이고 $N=2M$으로 짝수를 맞추면 실제 여유는 $30.0$과 $35.7$로 그보다 크다.

\begin{figure}[htbp]
\centering
\begin{tikzpicture}
\begin{semilogyaxis}[sdplot,xlabel={$N-\mathrm{top}$},ylabel={상대오차},xmin=0,xmax=58,ymin=1e-17,ymax=1,ytick={1e-16,1e-12,1e-8,1e-4,1}]
\addplot[black,mark=*,mark size=1.2pt] table[x=d,y=err]{figures/06/miller_n_x10_n0.dat};
\addlegendentry{$J_{0.3}(10)$}
\addplot[black,dashed,mark=square*,mark size=1.2pt] table[x=d,y=err]{figures/06/miller_n_x10_n20.dat};
\addlegendentry{$J_{20.3}(10)$}
\addplot[black!45,thin] coordinates {(30.0,1e-17) (30.0,1)};
\addlegendentry{stardust의 $N$}
\addplot[black!45,thin,forget plot] coordinates {(35.7,1e-17) (35.7,1)};
\end{semilogyaxis}
\end{tikzpicture}
\caption{Miller 방법(Neumann 합으로 정규화)의 오차를 출발 차수의 여유 $N-\mathrm{top}$에 따라 그렸다. 세로선은 stardust가 두 경우에 고르는 여유이다(\code{figures/06/bessel.py}).}
\label{fig:miller-n}
\end{figure}

\subsection{Normalizing Sums}

척도를 맞추는 관계식은 함수마다 다르다. 정수 차수의 Bessel 함수에는 생성 함수
\[
e^{(x/2)(t-1/t)}=\sum_{n\in\mathbb Z}J_n(x)t^n
\]
가 있다. 왼쪽을 $e^{xt/2}e^{-x/(2t)}$의 두 급수의 곱으로 펼쳐 $t^n$의 계수를 모으면 \cref{eq:jv-series}가 나오므로 생성 함수는 정의에서 바로 따라 나온다. $t=1$을 넣고 $J_{-n}=(-1)^nJ_n$을 쓰면
\[
1=J_0(x)+2\sum_{k\ge1}J_{2k}(x)
\]
이다. 이 생성 함수는 $t$에 대한 Laurent 급수(5장 \cref{thm:laurent})이므로 계수의 적분 표현을 $|t|=1$에서 쓰면 Bessel의 적분 $J_n(x)=\pi^{-1}\int_0^\pi\cos(n\theta-x\sin\theta)\,d\theta$가 나오고 여기서 $|J_n(x)|\le1$이다. 이 부등식 덕분에 정규화 합의 항들도 넘치지 않는다. \cref{fig:bessel-dir}의 Miller 열은 이것으로 정규화했다. 정수가 아닌 차수 $\nu=n+\mu$에는 $(x/2)^\mu$의 Neumann 전개
\[
\frac{(x/2)^\mu}{\Gamma(\mu+1)}=J_\mu(x)+\sum_{k\ge1}\frac{(\mu+2k)\,(\mu+1)_{k-1}}{k!}J_{\mu+2k}(x)
\]
를 쓴다(\cite{dlmf} 10.23.15). $\mu=0$이면 앞의 식이 된다. 증명에는 계수 비교와 망원합 하나면 된다.

\begin{proposition}
\label{prop:neumann}
$\mu>-1$, $x>0$이면 위 Neumann 전개가 성립한다.
\end{proposition}

\begin{proof}
오른쪽의 각 $J_{\mu+2k}$에 \cref{eq:jv-series}를 넣으면 이중급수가 절대수렴하므로 $(x/2)^{\mu+2m}$의 계수를 모을 수 있다. $k=0$의 계수를 $c_0=1$, $k\ge1$의 계수를 $c_k=(\mu+2k)\Gamma(\mu+k)/(\Gamma(\mu+1)k!)$로 쓰면($k=0$에도 같은 식이 $1$을 준다) $(x/2)^{\mu+2m}$의 계수는
\[
\frac{(-1)^m}{\Gamma(\mu+1)}S_m,\qquad S_m=\sum_{k=0}^m(-1)^k\frac{(\mu+2k)\Gamma(\mu+k)}{k!\,(m-k)!\,\Gamma(\mu+m+k+1)}
\]
이다. $m=0$이면 $S_0=\mu\Gamma(\mu)/\Gamma(\mu+1)=1$이고 이것이 왼쪽의 계수이다. $m\ge1$에서 $S_m=0$임을 보이면 된다. $W_k=(-1)^{k-1}\Gamma(\mu+k)/((k-1)!\,(m-k)!\,\Gamma(\mu+m+k))$($1\le k\le m$), $W_0=W_{m+1}=0$으로 두면
\[
W_{k+1}-W_k=\frac{(-1)^k\Gamma(\mu+k)}{k!\,(m-k)!\,\Gamma(\mu+m+k+1)}\bigl[(\mu+k)(m-k)+k(\mu+m+k)\bigr]
\]
이고 괄호 안은 $m(\mu+2k)$이다. 그래서 $S_m$의 각 항은 $(W_{k+1}-W_k)/m$이고 $S_m=(W_{m+1}-W_0)/m=0$이다.
\end{proof}

\noindent 이 증명은 \code{figures/06/neumann.py}에서 정확한 유리수로도 확인했다. 차수가 $x$보다 큰 항은 모두 양수이지만 $x$보다 작은 차수의 항은 진동하므로 합에는 상쇄가 있다. 그 정도를 조건수 $\kappa=\sum|\text{항}|/|\text{합}|$으로 재면 $\mu=0.3$에서 $x=10$이면 $2.0$, $x=30$이면 $2.9$, $x=100$이면 $5.2$이고 $\mu=-0.4$에서는 $x=100$에서 $10.7$이다(\cref{fig:neumann}). 정규화는 반올림을 이 배수만큼만 키우고 그 배수는 $x$와 함께 천천히 자란다. stardust의 \code{fast\_j}는 $\nu<500$(\code{kFastOrder})이고 $x\ge2^{-200}$이며 앞 절의 다른 경로(급수, Hankel, 닫힌 꼴의 위쪽 점화)가 맞지 않을 때 이 경로를 쓴다(\cref{alg:miller-bessel}, 경로 전체는 \cref{tab:bessel-paths}).

\begin{figure}[htbp]
\centering
\begin{tikzpicture}
\begin{semilogxaxis}[sdplot,xlabel={$x$},ylabel={조건수 $\kappa$},xmin=0.4,xmax=120,ymin=0,ymax=12]
\addplot[black,mark=*,mark size=1.2pt] table[x=x,y=k03]{figures/06/neumann.dat};
\addlegendentry{$\mu=0.3$}
\addplot[black,dashed,mark=square*,mark size=1.2pt] table[x=x,y=km04]{figures/06/neumann.dat};
\addlegendentry{$\mu=-0.4$}
\end{semilogxaxis}
\end{tikzpicture}
\caption{Neumann 합의 조건수 $\sum|\text{항}|/|\text{합}|$(십진 산술; \code{figures/06/neumann.py}).}
\label{fig:neumann}
\end{figure}

\begin{algorithm}[htbp]
\begin{algorithmic}[1]
\Require $\nu=n+\mu$ ($|\mu|\le1/2$), $x>0$
\Ensure $J_\nu(x)$
\State $\mathrm{top}\gets\max(\nu,x)$; $M\gets\lfloor(\mathrm{top}+10\,\mathrm{top}^{1/3}+8)/2\rfloor+1$; $N\gets2M$
\State $(1/x)_h+(1/x)_l\gets$ double-double $1/x$ \Comment{계수 $2v/x$를 한 번만 반올림}
\State $f_{N+1}\gets0$, $f_N\gets1$; $r\gets1$; $S\gets0$
\For{$m=M$ \textbf{downto} $1$} \Comment{한 번에 두 차수: 짝수 차수에서 합을 누적}
  \State $S\gets\mathrm{fma}((\mu+2m)r,\ f_{2m},\ S)$; $r\gets r\cdot m/(\mu+m-1)$ ($m\ge2$) \Comment{$r=b_m/b_M$}
  \State $f_{2m-1}\gets\mathrm{fma}(c_{2m},f_{2m},-f_{2m+1})$; $f_{2m-2}\gets\mathrm{fma}(c_{2m-1},f_{2m-1},-f_{2m})$ \Comment{$c_k=2(\mu+k)/x$}
  \State $f_n$을 지나면 기억한다; $|f_{2m-2}|>2^{600}$이면 모두 $2^{-600}$배
\EndFor
\State $\Sigma\gets f_0+S/r$ \Comment{Neumann 합, $f$와 같은 단위}
\State \Return $f_n\cdot(x/2)^\mu/(\Gamma(1+\mu)\,\Sigma)$
\end{algorithmic}
\caption{Neumann 합으로 정규화한 Miller 방법 (stardust \code{bessel.hpp}의 \code{miller}, \code{fast\_j})}
\label{alg:miller-bessel}
\end{algorithm}

\cref{alg:miller-bessel}의 각 단계를 보자. 1행의 $N$은 앞 소절의 여유이고 짝수로 맞춘다. 4--7행은 루프 한 번에 두 차수를 내려간다. 그러면 Neumann 합의 항이 짝수 차수에만 있다는 것을 분기 없이 처리할 수 있다. 계수 $b_m=(\mu+1)_{m-1}/m!$은 $m^{\mu-1}$처럼 천천히 준다. 위에서 내려가며 $b_{m-1}=b_m\,m/(\mu+m-1)$로 만들려면 시작값 $b_M$이 필요하다. 그 값 대신 $b_M$을 $1$로 둔 비 $r=b_m/b_M$을 계산하다가 끝에서 $r=b_1/b_M=1/b_M$이 되면 $S/r$로 척도를 되돌린다. 6행의 계수 $c_k=2(\mu+k)/x$에서 $k\le n$이면 $\mu+k=\nu-(n-k)$가 정확한 double이다. 그보다 큰 차수에서는 $\mu+k$가 반올림될 수 있지만 그 오차는 걸음마다의 반올림과 같은 크기이다. $1/x$는 double-double로 두어 계수를 \code{fma}로 한 번만 반올림한다. $1/x$를 binary64로 한 번 반올림하면 $x$ 자체가 1 ulp 움직이는 셈이고 결과가 $xJ'/J$ ulp만큼 움직이기 때문이다. 7행의 척도 맞춤은 $x$가 $N$에 비해 작아 $\tilde f$가 $2^{1024}$를 넘을 수 있는 경우를 위한 것이다. 같은 함수를 Airy 함수도 쓴다(\code{csrc/special/ops/airy.hpp}의 \code{negative\_bessel}). $-9.7<x<-2$에서 $\operatorname{Ai}$, $\operatorname{Bi}$와 도함수는 $\zeta=(2/3)|x|^{3/2}$에서의 $J_{\pm1/3}$, $J_{\pm2/3}$의 조합이다(\cite{dlmf} 9.6.6--9.6.9). 차수 $-1/3+k$와 $-2/3+k$는 두 격자를 이루므로 Miller 점화를 두 번 돌려 네 값을 모두 얻는다. $\zeta$는 double-double로 만들고 그 낮은 부분은 $J'_v=J_{v-1}-(v/\zeta)J_v$를 써서 일차 보정으로 더한다. $2$의 거듭제곱이므로 정확하고 비율을 바꾸지 않는다. 비용은 $N\approx\mathrm{top}+10\,\mathrm{top}^{1/3}$에 비례하는 \code{fma} 몇 번씩이다.

\subsection{Pincherle's Theorem and the Continued Fraction CF1}

최소해는 3장의 연분수(\cref{sec:continued-fractions})와 바로 이어진다. \cref{eq:three-term}을 비 $\rho_n=y_n/y_{n-1}$에 대해 정리하면 $\rho_n=-b_n/(a_n+\rho_{n+1})$이고 이 식을 되풀이하면
\[
\frac{y_n}{y_{n-1}}=\cfrac{-b_n}{a_n-\cfrac{b_{n+1}}{a_{n+1}-\cfrac{b_{n+2}}{a_{n+2}-\cdots}}}
\]
로 펼쳐진다. 연분수를 $N$층에서 자르는 것은 $\rho_{N+1}=0$, 곧 $\tilde f_{N+1}=0$으로 놓는 것과 같으므로 잘린 연분수는 바로 \cref{alg:miller}의 비이다. \cref{thm:miller}의 증명이 곧 다음 정리의 한쪽 방향인 셈이다.

\begin{theorem}[Pincherle]
\label{thm:pincherle}
\cref{eq:three-term}의 위 연분수가 수렴할 필요충분조건은 점화식이 최소해 $f$를 갖는 것이다. 수렴하면 그 값은 $f_n/f_{n-1}$이다.
\end{theorem}

최소해가 있으면 잘린 연분수가 $\tilde f^{(N)}_n/\tilde f^{(N)}_{n-1}$이고 \cref{thm:miller}의 증명에서 이것이 $f_n/f_{n-1}$로 간다($f_{n-1}\ne0$이면). 역방향을 보이자. 위치 $n$의 연분수를 $N$층에서 자른 근사분수를 $A_N/B_N$이라 하면 3장의 근사분수 점화식에서 $A_N$, $B_N$은 $N$에 대해 \cref{eq:three-term}과 같은 꼴의 점화식을 만족한다. 원래 점화식의 두 해에서 첨자만 옮긴 꼴이다. 둘의 Casorati 행렬식은 $\prod(-b_j)$꼴로 $0$이 아니므로 독립이다. 연분수가 $L$로 수렴하면 $h=A-LB$도 해이고 $h_N/B_N=A_N/B_N-L\to0$이다. 다른 해 $g=\alpha A+\beta B$를 아무것이나 잡으면 $h_N/g_N=(A_N/B_N-L)/(\alpha A_N/B_N+\beta)$이고 분모는 $\alpha L+\beta$로 간다. 그러므로 $\alpha L+\beta\ne0$이면($g$가 $h$의 상수배가 아니면) $h_N/g_N\to0$이고 $h$는 최소해이다. 그 값이 비 $f_n/f_{n-1}$임은 앞 방향의 논증에서 나온다\cite{gautschi1967recurrence}. 최소해의 이웃한 두 값의 비는 3장의 Lentz 알고리즘(\cref{alg:lentz})으로 앞에서부터 계산할 수 있다. 비 하나와 Wronskian 같은 관계식이 있으면 함수값 자체도 얻는다. Bessel 함수 계산에서는 이 연분수를 흔히 CF1이라 부른다. $J'_\nu/J_\nu=\nu/x-J_{\nu+1}/J_\nu$이고 \cref{eq:bessel-rec}에서 $a_k=-2(\nu+k)/x$, $b_k=1$이므로
\[
\frac{J'_\nu}{J_\nu}=\frac\nu x-\cfrac{1}{2(\nu+1)/x-\cfrac{1}{2(\nu+2)/x-\cdots}}
\]
이다.

stardust의 일반 경로(\code{jy\_temme}, Numerical Recipes의 \code{bessjy}\cite{press2007nr}를 고친 것)는 CF1을 두 번 지난다(\cref{alg:cf1}). 앞으로 가는 Lentz 반복은 수렴할 때까지 필요한 항의 수만 센다. 그다음 같은 수의 항을 뒤에서부터 다시 계산해 값을 얻는다. 3장에서 보았듯 Lentz의 앞으로 가는 곱은 걸음마다의 반올림을 곱해 쌓지만 뒤에서부터 계산하면 각 걸음의 반올림이 감쇠된다(\cref{prop:cf-rounding}, \cref{alg:lentz-backward}). 변형 Bessel 함수의 CF1 $I'_\nu/I_\nu=\nu/x+1/(2(\nu+1)/x+1/(2(\nu+2)/x+\cdots))$은 항이 모두 양수이므로 3장의 \code{lentz\_backward}를 그대로 쓴다. 근사분수의 분자와 분모가 점화식을 만족한다는 사실은 3장의 \code{cf\_fast}(\cref{alg:cf-fast})가 직접 쓴다. 그 함수는 Lentz의 나눗셈 대신 $A_n=b_nA_{n-1}+a_nA_{n-2}$, $B_n=b_nB_{n-1}+a_nB_{n-2}$를 \code{fma} 두 번씩으로 돌리고 $2^{\pm256}$의 정확한 척도로 크기를 묶는다.

\begin{algorithm}[htbp]
\begin{algorithmic}[1]
\Require $\nu\ge0$, $x>0$
\Ensure $(J_\nu,J'_\nu)$에 공통인 양의 인자를 곱한 쌍
\State $T\gets$ Lentz 앞 반복이 $|\Delta-1|\le\varepsilon$에 닿은 항의 수 \Comment{값은 버린다}
\State $d\gets\infty$; $c\gets0$
\For{$k=T+4$ \textbf{downto} $2$}
  \State $d\gets2(\nu+k)/x-1/d$ \Comment{$d_k=J_{\nu+k-1}/J_{\nu+k}$}
  \State \textbf{if} $\nu+k<x+2$ \textbf{and} $d<0$ \textbf{then} $c\gets c+1$ \Comment{부호 바뀜을 센다}
\EndFor
\State $s\gets(-1)^c$ \Comment{$J_{\nu+1}$의 부호}
\State $d\gets2(\nu+1)/x-1/d$ \Comment{$J_\nu/J_{\nu+1}$}
\State \textbf{if} $|d|\le2^{500}$ \textbf{then} \Return $(sd,\ (\nu/x)sd-s)$ \Comment{$|J_{\nu+1}|$ 단위}
\State \Return $(\operatorname{sgn}(d)s,\ (\nu/x-1/d)\operatorname{sgn}(d)s)$ \Comment{$J_{\nu+1}\approx0$: $|J_\nu|$ 단위}
\end{algorithmic}
\caption{CF1의 뒤 패스와 부호 세기 (stardust \code{bessel.hpp}의 \code{jy\_temme})}
\label{alg:cf1}
\end{algorithm}

\cref{alg:cf1}에는 두 가지 세부가 있다. 먼저 부호를 보자. 큰 차수의 $J$는 양수이므로 $d_k=J_{\nu+k-1}/J_{\nu+k}$의 부호가 바뀔 때마다 $J$의 부호가 바뀐다. 음수인 $d_k$의 수를 세면 $J_{\nu+1}$의 부호가 나온다. 다음 보조정리 때문에 $\nu+k<x+2$인 차수만 세면 된다.

\begin{lemma}
\label{lem:cf1-positive}
뒤 패스의 $d_k=2(\nu+k)/x-1/d_{k+1}$을 $d_{T+5}=\infty$에서 시작하면 $\nu+k\ge x+2$인 모든 $k$에서 $d_k>1$이다.
\end{lemma}

\begin{proof}
$d_{k+1}>1$이면 $d_k>2(\nu+k)/x-1\ge2(x+2)/x-1>1$이다. 시작값 $\infty$는 $1$보다 크고 $\nu+k$는 $k$가 줄수록 작아지므로 조건이 성립하는 동안 귀납이 이어진다.
\end{proof}

\noindent 그래서 코드는 $\nu+k\ge x+2$인 동안은 부호를 세지 않고 값만 돈다. 그 아래는 부호가 예측할 수 없이 바뀌므로 분기 없이 $d<0$의 진릿값을 더해 센다. 앞 패스에서 부호를 얻으면 반올림 때문에 $J_\nu$의 영점의 반대쪽에 설 수 있고 그러면 $J'_\nu$와 $Y_\nu$의 부호가 뒤집힌다. 둘째로 이 알고리즘은 비 자체를 만들지 않는다. $x$가 $J_\nu$의 영점에 가장 가까운 double이면 $d_1$이 정확히 $0$이 될 수 있다. 그때 $J'/J$는 무한대이고 그것과의 곱은 모두 NaN이 된다. 이를 피하려고 쌍 $(J_\nu,J'_\nu)$을 공통 인자를 곱한 채로 들고 다니고 그 인자는 마지막에 Wronskian으로 정한다(다음 절). 항의 수 $T$는 $x>\nu$이면 대략 $x$에 비례하고 $x<\nu$이면 작다(\cref{fig:cf1-terms}). $x>\nu$에서 연분수의 꼬리 $2(\nu+k)/x$가 $2$ 아래인 동안은 근사분수가 진동하는 중립 영역이기 때문이다. 반복의 상한은 \code{kMaxCfTerms}$=10^7$이고 그 안에 수렴하지 않으면 NaN을 낸다. $x\ge2^{53}$이고 $\nu-x<10^{-12}x$이면 걸음 $2/x$가 $2\nu/x$의 반 ulp보다 작아 계수가 사실상 상수 $2+\delta$인 연분수가 되고 수렴에 $18/\sqrt\delta$ 항쯤이 필요하므로 미리 NaN을 낸다. 중립 영역의 CF1은 반올림도 쌓는다. 코드의 주석에 따르면 $x>\nu$에서 CF1은 약 $0.3x$ ulp를 잃으므로 $x\ge30$이고 $\nu<x$이면 stardust는 CF1 대신 작은 차수에서 Hankel 전개로 씨앗을 만들고 double-double로 위로 점화한다(\code{jy\_upward}).

\begin{figure}[htbp]
\centering
\begin{tikzpicture}
\begin{loglogaxis}[sdplot,xlabel={$x$},ylabel={CF1의 항 수 $T$},xmin=0.1,xmax=3000,ymin=1,ymax=5000]
\addplot[black,mark=*,mark size=1.0pt] table[x=x,y=t0]{figures/06/cf1terms.dat};
\addlegendentry{$\nu=0.3$}
\addplot[black,dashed,mark=square*,mark size=1.0pt] table[x=x,y=t10]{figures/06/cf1terms.dat};
\addlegendentry{$\nu=10.3$}
\addplot[black!50,mark=triangle*,mark size=1.2pt] table[x=x,y=t100]{figures/06/cf1terms.dat};
\addlegendentry{$\nu=100.3$}
\end{loglogaxis}
\end{tikzpicture}
\caption{Lentz 앞 반복이 $|\Delta-1|\le2^{-52}$에 닿는 CF1의 항 수(\code{figures/06/cf1terms.py}).}
\label{fig:cf1-terms}
\end{figure}


\section{Forward Recurrence for Dominant Solutions}
\label{sec:dominant}

\subsection{Two Paths in stardust}

지배해 $Y_\nu$와 $K_\nu$는 위로 점화하면 안정하다(\cref{prop:forward-error}). 그러니 낮은 두 차수의 값, 곧 씨앗만 얻으면 된다. stardust의 Bessel 함수는 차수와 인자에 따라 길을 나눈다(\cref{tab:bessel-paths}). 차수 $|\nu|<500$(\code{kFastOrder})이고 $x\ge2^{-200}$이면 먼저 binary64만 쓰는 빠른 경로(\code{fast\_j}, \code{fast\_y}, \code{fast\_ik\_scaled})를 시도한다. 차수가 더 크거나 빠른 경로의 결과가 정규수가 아니거나 루프가 끝나지 않으면 double-double 점화를 쓰는 일반 경로(\code{jy\_temme}, \code{ik\_temme})로 간다. 코드의 주석에 따르면 차수 $500$ 아래에서 빠른 경로의 오차는 참값 파일 기준 포락선 $13$·$15$ ulp 이하로 일반 경로와 같은 수준이고 세 배에서 네 배 빠르다. 그보다 큰 차수에서는 binary64 점화가 약 $n/30$ ulp를 잃으므로 일반 경로로 넘긴다.

\begin{table}[htbp]
\begin{sdtable}{colspec={lX}, row{1}={font=\small\bfseries}}
\toprule
경우 & 방법 \\
\midrule
\code{jv}, 빠른 경로 & $x^2<2(\nu+1)$: 급수 \cref{eq:jv-series}; $x\ge25$이고 받아들여지면 Hankel; 차수 $0,1$: 유리 근사; $\nu<x$: 씨앗($\mu=-1/2$의 닫힌 꼴, $x\ge25$의 Hankel)에서 위로 점화; 그 밖: Miller와 Neumann 합(\cref{alg:miller-bessel}) \\
\code{yv}, 빠른 경로 & Hankel; 아니면 씨앗 $Y_\mu,Y_{\mu+1}$(닫힌 꼴, Hankel, $x<3$이면 Temme 급수, 그 밖은 CF2와 Miller의 $J$)에서 위로 점화(\cref{alg:fast-y}) \\
\code{iv}, \code{kv}, 빠른 경로 ($|\nu|<30$, $x<700$) & $x^2<4(\nu+1)$이면 $I$의 급수; 아니면 $K_\mu,K_{\mu+1}$(Temme, CF2, 큰 $x$ 급수)에서 $K$를 위로, $I$는 Wronskian과 역방향 비(\cref{alg:ik}) \\
\code{iv} ($|\nu|<30$), \code{kv} ($|\nu|<400$) & 빠른 경로가 실패하면 Boost 경로; $|\nu|\ge30$의 $I,K$는 Debye 균등 전개(92장) \\
일반 경로 $J,Y$ & $x\ge30$: Hankel, 아니면 $\nu<x$에서 Hankel 씨앗과 double-double 위쪽 점화(\code{jy\_upward}); 그 밖: \cref{alg:jy} \\
\bottomrule
\end{sdtable}
\caption{stardust의 실수 차수 Bessel 함수가 고르는 길(\code{csrc/special/ops/bessel.hpp}; 음의 차수는 반사 공식, 영역의 근사식은 91--92장).}
\label{tab:bessel-paths}
\end{table}

\subsection{Seeds from Temme's Series and Steed's CF2}

씨앗은 Temme의 방법\cite{temme1975bessel,press2007nr}으로 얻는다. $x$가 작으면 \cref{prop:temme-series}의 급수로 $Y_\mu$, $Y_{\mu+1}$을 직접 계산한다(\cref{alg:temme}). 변형 Bessel 함수에도 같은 꼴의 급수가 있어 $K_\mu$, $K_{\mu+1}$을 준다.

\begin{algorithm}[htbp]
\begin{algorithmic}[1]
\Require $|\mu|\le1/2$, $0<x<3$
\Ensure $Y_\mu(x)$, $Y_{\mu+1}(x)$
\State $\Gamma_1,\Gamma_2,1/\Gamma(1\pm\mu)\gets$ $1/\Gamma(1+z)$의 Taylor 급수에서 \Comment{$\mu\to0$에서도 매끄럽게}
\State $d\gets\ln(2/x)$; $\sigma\gets\mu d$; $f\gets$ \cref{prop:temme-series}의 $f_0$ \Comment{$\sinh\sigma/\sigma$, $\mu\pi/\sin\mu\pi$는 $0$ 근처에서 $1$}
\State $p\gets(x/2)^{-\mu}\Gamma(1+\mu)/\pi$; $q\gets(x/2)^\mu\Gamma(1-\mu)/\pi$; $r\gets(2/\mu)\sin^2(\mu\pi/2)$; $c\gets1$
\State $S\gets f+rq$; $S_1\gets p$
\For{$k=1,2,\dots$}
  \State $f\gets(kf+p+q)/(k^2-\mu^2)$; $c\gets c\cdot(-x^2/4)/k$; $p\gets p/(k-\mu)$; $q\gets q/(k+\mu)$
  \State $\Delta\gets c(f+rq)$; $S\gets S+\Delta$; $S_1\gets S_1+c\,p-k\Delta$
  \State \textbf{if} $|\Delta|\le2^{-56}(1+|S|)$ \textbf{then} \Return $(-S,\ -2S_1/x)$
\EndFor
\end{algorithmic}
\caption{Temme의 급수 (stardust \code{bessel.hpp}의 \code{temme\_y})}
\label{alg:temme}
\end{algorithm}

\cref{alg:temme}의 7행은 $\Delta=c_kg_k$를 다시 써서 $S_1$에 $c_k(p_k-kg_k)$를 더한다. $|c_k|=(x/2)^{2k}/k!$이고 $g_k$도 대략 $1/k!$처럼 줄므로 항은 $(x/2)^{2k}/(k!)^2$처럼 줄어 $x<3$에서 수십 항이면 끝난다. $x$가 커지면 이 항들이 커졌다가 줄어 상쇄가 생긴다.

$x$가 크면 Steed의 연분수 CF2\cite{barnett1974}를 쓴다. Hankel 함수 $H^{(1)}_\mu=J_\mu+iY_\mu$의 로그 도함수
\[
p+iq=\frac{J'_\mu+iY'_\mu}{J_\mu+iY_\mu}=-\frac{1}{2x}+i+\frac ix\cdot\cfrac{(1/2)^2-\mu^2}{2(x+i)+\cfrac{(3/2)^2-\mu^2}{2(x+2i)+\cdots}}
\]
는 $H^{(1)}_\mu(x)=\sqrt{2/(\pi x)}e^{i\chi}U(\cdot)$의 꼴에서 합류 초기하 함수 $U$의 매개변수에 대한 삼항 점화식에 최소해가 있으므로 \cref{thm:pincherle}에 따라 연분수로 나타난다. 유도는 Barnett 등\cite{barnett1974}과 Numerical Recipes\cite{press2007nr}(§6.6)에 있다. 이 연분수는 $x$가 클수록 빨리 수렴한다. 코드의 주석에는 $x=2$에서 $40$항, $x=3$에서 $30$항쯤으로 적혀 있다. 두 방법을 바꾸는 경계는 경로마다 다르다. 일반 경로는 Numerical Recipes처럼 $x<2$에서 Temme 급수를 쓰고 빠른 경로의 \code{fast\_y}는 $x<3$(\code{kTemmeBelow})에서 쓴다. 같은 주석에 적힌 측정으로는 $\nu<3$, $x\in[1.5,6]$의 Arb 참값 묶음에서 경계를 $2$, $3$, $4$에 두면 최대 오차가 포락선 $32$, $21$, $35$ ulp였다.

CF2는 비 $p+iq$만 주므로 척도는 Wronskian으로 정한다. CF1(또는 Miller)과 아래쪽 점화로 얻은 $J_\mu$, $J'_\mu$에는 공통 인자 $c$가 붙어 있다. CF2를 쓰면 $J'=pJ-qY$, $Y'=qJ+pY$이므로 Wronskian $J_\mu Y'_\mu-J'_\mu Y_\mu=2/(\pi x)$은 $q(J^2+Y^2)$이 된다. 이 관계를 쓰면 $Y_\mu=c\,(pj_a-j_b)/q$, $c=\sqrt{2/(\pi xq(j_a^2+y_r^2))}$($y_r=(pj_a-j_b)/q$)처럼 $J_\mu$로 나누지 않고 정해진다($j_a,j_b$는 큰 쪽이 $1$이 되게 척도를 맞춘 $J_\mu/c$, $J'_\mu/c$). $J_\mu$로 나누면 안 되는 까닭은 $\mu=-1/2$에서 $J_\mu=\sqrt{2/(\pi x)}\cos x$가 $x=\pi/2$에서 $0$이 되기 때문이다.

일반 경로(\cref{alg:jy})에서 $\mu$를 정하는 방법은 $x$에 따라 다르다. $x<2$이면 $n_l=\operatorname{round}(\nu)$로 두어 $|\mu|\le1/2$이다. $x\ge2$이면 $n_l=\max(0,\lfloor\nu-x+1.5\rfloor)$로 두어 $\mu$가 $x$ 근처까지 커질 수 있다. CF2는 $x\ge2$에서 $\mu$와 상관없이 수렴하므로 큰 $\mu$에서 계산하고 점화는 $\nu-x$걸음만 한다. 그러면 점화가 중립인 $n<x$ 영역을 걷지 않는다.

\begin{algorithm}[htbp]
\begin{algorithmic}[1]
\Require $\nu\ge0$, $x>0$; $\nu=n_l+\mu$ (위의 $n_l$)
\Ensure $J_\nu,Y_\nu,J'_\nu,Y'_\nu$ (각각 $m\cdot2^e$ 꼴)
\State $(a,b)\gets$ \cref{alg:cf1} \Comment{$(J_\nu,J'_\nu)/c$}
\State $e_J\gets0$; \textbf{for} $l=n_l-1$ \textbf{downto} $0$: double-double로 $(a,b)$를 한 차수 내린다; $|a|$가 크면 $2^{-k}$배하고 $e_J\mathrel{+}=k$
\State $(j_a,j_b)\gets(a,b)/\max(|a|,|b|)$ \Comment{$(J_\mu,J'_\mu)/c'$}
\State $x<2$: Temme 급수(\cref{alg:temme})로 $Y_\mu,Y_{\mu+1}$, $c'\gets(2/(\pi x))/(j_aY'_\mu-j_bY_\mu)$; $x\ge2$: CF2와 Wronskian으로 $c'$, $Y_\mu$, $Y'_\mu$
\State $J_\nu\gets(\text{1행의 }a)\cdot c'\cdot2^{-e_J}\cdots$ \Comment{척도를 되돌린다}
\State $e_Y\gets0$; \textbf{for} $i=1$ \textbf{to} $n_l$: double-double로 $Y_{\mu+i+1}=(2(\mu+i)/x)Y_{\mu+i}-Y_{\mu+i-1}$; 크면 $2^{-k}$배, $e_Y\mathrel{+}=k$
\State \Return $J_\nu$, $Y_\nu\cdot2^{e_Y}$, $J'_\nu$, $Y'_\nu=(\nu/x)Y_\nu-Y_{\nu+1}$
\end{algorithmic}
\caption{Temme--Steed 방법의 골격 (stardust \code{bessel.hpp}의 \code{jy\_temme})}
\label{alg:jy}
\end{algorithm}

\cref{alg:jy}의 비용은 CF1의 $T+4$항(\cref{fig:cf1-terms}), 아래쪽 점화 $n_l$걸음, Temme 급수나 CF2의 수십 항, 위쪽 점화 $n_l$걸음이다. 점화 한 걸음은 double-double 곱 두 번과 덧셈 한 번이다. 빠른 경로의 \code{fast\_y}는 같은 생각을 binary64로 한다(\cref{alg:fast-y}). $J_\mu$, $J'_\mu$를 CF1 대신 짧은 Miller 점화(\code{miller<false>}, 정규화 합 없이 $N=\mathrm{top}+8\,\mathrm{top}^{1/3}+6$, $\mathrm{top}=\max(x,1)$)로 얻고 공통 인자는 Wronskian으로 정한다.

\begin{algorithm}[htbp]
\begin{algorithmic}[1]
\Require $0\le\nu<500$, $x\ge2^{-200}$; $\nu=n+\mu$, $|\mu|\le1/2$
\Ensure $Y_\nu(x)$, 또는 ``일반 경로로''
\State \textbf{if} $x\ge25$ \textbf{and} Hankel 전개가 $\nu$에서 받아들여지면 \Return 그 값
\If{$\mu=0$ \textbf{or} $\mu=-1/2$ \textbf{or} $x\ge25$} $(Y_\mu,Y_{\mu+1})\gets$ 유리 근사, 닫힌 꼴, 또는 Hankel
\ElsIf{$x<3$} $(Y_\mu,Y_{\mu+1})\gets$ \cref{alg:temme}
\Else
  \State $p+iq\gets$ CF2; $(f_0,f_1)\gets$ Miller의 비(정규화 없이) \Comment{$(J_\mu,J_{\mu+1})/c$}
  \State $(j_a,j_b)\gets(f_0,(\mu/x)f_0-f_1)/\max$; $y_r\gets(pj_a-j_b)/q$; $c\gets\sqrt{2/(\pi xq(j_a^2+y_r^2))}$
  \State $Y_\mu\gets cy_r$; $Y_{\mu+1}\gets(\mu/x)Y_\mu-(q\,c\,j_a+pY_\mu)$ \Comment{$Y'=qJ+pY$}
\EndIf
\State \Return \code{upward}$(\mu,x,n,Y_\mu,Y_{\mu+1})$ \Comment{\code{fma}와 double-double $1/x$의 계수}
\end{algorithmic}
\caption{$Y_\nu$의 빠른 경로 (stardust \code{bessel.hpp}의 \code{fast\_y})}
\label{alg:fast-y}
\end{algorithm}

변형 Bessel 함수의 빠른 경로(\cref{alg:ik})는 $K$가 지배해이고 $I$가 최소해라는 구조를 그대로 쓴다. $K_\mu$, $K_{\mu+1}$은 $x<2$에서 Temme 급수, $x\ge20$(\code{kKAsymFrom})에서 큰 $x$ 급수, 그 사이에서 Temme 꼴의 CF2로 얻는다. 큰 $x$ 급수의 가장 작은 항이 $e^{-2x}$ 정도이므로 $x=20$부터는 그 항이 $2^{-57}$ 아래이다. $K$의 CF2는 Numerical Recipes의 꼴을 그대로 쓰면 항마다 나눗셈 $d_i=1/(b_i+a_id_{i-1})$이 의존 사슬에 놓여 한 항에 $18$ 사이클쯤 든다. stardust는 $d_i=D_{i-1}/D_i$, $D_i=b_iD_{i-1}+a_iD_{i-2}$로 바꾸어(다시 근사분수의 분모 점화식이다) 사슬에는 \code{fma}와 곱만 남기고 나눗셈은 사슬 밖으로 뺀다. $D$와 그에 딸린 곱 $P$는 $2^{\pm400}$ 안에 머물도록 함께 척도를 맞춘다. $K$를 위로 점화한 뒤 $I$는 Wronskian $I_\nu K_{\nu+1}+I_{\nu+1}K_\nu=1/x$와 역방향 점화로 얻은 비 $I_{\nu+1}/I_\nu$에서 나온다.

\begin{algorithm}[htbp]
\begin{algorithmic}[1]
\Require $0\le\nu<30$, $0<x<700$; $\nu=n+\mu$, $|\mu|\le1/2$
\Ensure $e^xK_\nu(x)$, $e^{-x}I_\nu(x)$, 또는 ``다른 경로로''
\State $(k_0,k_1)\gets e^x(K_\mu,K_{\mu+1})$ \Comment{Temme ($x<2$), 큰 $x$ 급수 ($x\ge20$), CF2}
\For{$k=1,\dots,n$} $(k_0,k_1)\gets(k_1,\ \mathrm{fma}(c_{\mu+k},k_1,k_0))$ \Comment{$K_{v+1}=(2v/x)K_v+K_{v-1}$}
\EndFor
\State $N\gets\lfloor\sqrt{40x}\rfloor+10$; $f_{N+1}\gets0$, $f_N\gets1$
\For{$k=N$ \textbf{downto} $1$} $f_{k-1}\gets\mathrm{fma}(c_{\nu+k},f_k,f_{k+1})$; 크면 $2^{-600}$배
\EndFor \Comment{$I_{\nu+k}$의 비: 모든 항이 양수}
\State $\rho\gets f_1/f_0$ \Comment{$I_{\nu+1}/I_\nu$}
\State \Return $k_0$, $1/(x(k_1+\rho k_0))$ \Comment{Wronskian; 둘 다 정규수인지 확인}
\end{algorithmic}
\caption{$K_\nu$와 $I_\nu$의 빠른 경로 (stardust \code{bessel.hpp}의 \code{fast\_ik\_scaled}, \code{i\_ratio})}
\label{alg:ik}
\end{algorithm}

\cref{alg:ik}의 역방향 점화에는 빼기가 없으므로 어느 걸음에서도 상쇄가 일어나지 않는다. 코드의 주석에는 $x=0.5$--$100$, $\nu\in[-1/2,1/2]$에서 잰 필요한 여유가 $\sqrt{30x}+4$로 적혀 있고 시작 차수의 여유 $\sqrt{40x}+10$은 그보다 넉넉하다. $I$의 Wronskian은 $K$의 척도($e^x$)와 $I$의 척도($e^{-x}$)가 서로 지워지므로 척도 맞춘 값끼리 바로 쓸 수 있다. 결과가 정규수가 아니거나 루프가 끝나지 않으면 Boost 경로로 넘긴다.

\subsection{Long Recurrences}

\cref{prop:discrete-green}에 따르면 안정한 방향의 점화에서도 걸음마다의 반올림은 쌓인다. 걸음 수 $n$이 수백이면 binary64의 $u$가 수십 배로 커진다. \cref{fig:threeway}에서는 세 가지 산술로 같은 점화를 400걸음 돌렸다. 첫째는 계수 $2v/x$를 나눗셈으로 반올림하고 곱과 뺄셈을 따로 반올림하는 보통의 binary64이다. 둘째는 stardust의 빠른 경로가 쓰는 식이다. $1/x$를 double-double로 두어 계수를 \code{fma}로 한 번만 반올림하고 한 걸음을 \code{fma} 하나로 끝낸다. 셋째는 일반 경로의 double-double이다. 참값은 십진 산술 $90$자리의 같은 점화이다. 아래로 점화한 $J$는 binary64로 최대 $40u$, 빠른 경로 식으로 $32u$, double-double로 $1.1u$였다. 위로 점화한 $Y$는 각각 $31u$, $43u$, $1.3u$였다. 오차는 짝의 크기 $\max(|y_v|,|y_{v\pm1}|)$에 대한 것이다. \code{fma}와 정확한 계수로 걸음마다의 반올림은 줄어도 쌓이는 양은 걸음 수에 따라 커진다. 그래서 stardust는 빠른 경로를 차수 $500$ 아래에서만 쓴다.

\begin{figure}[htbp]
\centering
\begin{tikzpicture}
\begin{groupplot}[group style={group size=1 by 2,vertical sep=1.5cm},sdplot,width=0.6\textwidth,height=0.3\textwidth,xmin=0,xmax=400,ymin=0.01,ymax=100,ytick={0.01,0.1,1,10,100},ymode=log,xlabel={차수 $v-0.3$},ylabel={오차$/u$}]
\nextgroupplot[title={\scriptsize $J$, 아래로}]
\addplot[black,only marks,mark=*,mark size=0.6pt] table[x=k,y=Jplain]{figures/06/threeway.dat};
\addlegendentry{binary64}
\addplot[black!45,only marks,mark=square*,mark size=0.6pt] table[x=k,y=Jfast]{figures/06/threeway.dat};
\addlegendentry{빠른 경로 식}
\addplot[black,only marks,mark=triangle*,mark size=0.8pt] table[x=k,y=Jdd]{figures/06/threeway.dat};
\addlegendentry{double-double}
\nextgroupplot[title={\scriptsize $Y$, 위로}]
\addplot[black,only marks,mark=*,mark size=0.6pt] table[x=k,y=Yplain]{figures/06/threeway.dat};
\addlegendentry{binary64}
\addplot[black!45,only marks,mark=square*,mark size=0.6pt] table[x=k,y=Yfast]{figures/06/threeway.dat};
\addlegendentry{빠른 경로 식}
\addplot[black,only marks,mark=triangle*,mark size=0.8pt] table[x=k,y=Ydd]{figures/06/threeway.dat};
\addlegendentry{double-double}
\end{groupplot}
\end{tikzpicture}
\caption{$x=100$에서 $J$를 차수 $400.3$에서 $0.3$으로, $Y$를 $0.3$에서 $400.3$으로 점화할 때 짝의 크기에 대한 오차(\code{figures/06/threeway.py}; $0.01$은 그 이하).}
\label{fig:threeway}
\end{figure}

계수의 차수도 정확해야 한다. 같은 실험을 차수 $\nu=286.3$에서 $0.3$까지 하면서(\code{figures/06/downrec.py}) 차수를 $\mu+k$로 계산해 반올림하면 걸음마다 계수가 틀려 double-double로도 $100u$를 잃었다. stardust처럼 차수를 $v=\nu-k$의 정확한 정수 뺄셈으로 만들면 $1.03u$였다.

넘침도 문제이다. $\nu\gg x$이면 정규화 전의 $J$가 차수를 내려갈수록 대략 $(2\nu/(ex))^\nu$배로 자라 binary64의 범위를 넘는다. stardust는 값을 $m\cdot2^e$ 꼴로 들고 다니며 $|m|$이 커지면 $2$의 거듭제곱으로 나누고 지수를 따로 센다(\code{rescale\_down}). 최종 결과도 $m\cdot2^e$ 꼴(\code{Scaled})이고 $J_\nu$가 $2^{-1074}$ 아래로 내려가는지는 마지막에 한 번만 판정한다. 이렇게 하면 중간값이 넘치거나 아래로 넘쳐도 결과를 잃지 않는다. 반복의 길이에는 상한(\code{kMaxOrderSteps})을 두어 한 번의 호출이 약 $0.25$초를 넘지 않게 한다. 그 이상이 필요하면 NaN을 낸다.

\subsection{Checking with the Wronskian}

Wronskian은 검증에도 쓴다. 정확한 $J$와 $Y$는 $J_{\nu+1}Y_\nu-J_\nu Y_{\nu+1}=2/(\pi x)$를 만족하므로 따로 계산한 네 값으로 이 식의 잔차를 재면 각 값의 오차를 대략 볼 수 있다. 따로 빌드한 stardust에서 $\nu\in[0,40]$($2^{-10}$의 배수), $x\in[0.5,30]$의 무작위 점 $300$개를 재면 잔차 $(\pi x/2)(J_{\nu+1}Y_\nu-J_\nu Y_{\nu+1})-1$을 $(\pi x/2)(|J_{\nu+1}Y_\nu|+|J_\nu Y_{\nu+1}|)$로 나눈 값이 최대 $1.2\times10^{-15}$이었다. 같은 점에서 \code{jv}를 \cref{eq:jv-series}의 십진 산술 $90$자리 합과 견주면 $|\hat J-J|/\sqrt{J^2+Y^2}$가 최대 $9.0\times10^{-16}$이었다. 상대오차로는 최대 $1.3\times10^{-14}$인데 그런 점은 모두 $J_\nu$의 영점 근처이다. 영점 근처에서 상대오차가 큰 것은 조건수 탓이므로 진동하는 함수의 정확도는 포락선 $\sqrt{J^2+Y^2}$에 대해 재야 공정하다.


\section{Recurrences as Boundary Value Problems}
\label{sec:olver}

\subsection{Olver's Method}

Miller의 방법은 동차 점화식의 최소해를 위한 것이다. 원하는 수열이 비동차 점화식
\begin{equation}
\label{eq:inhom}
c_ny_{n-1}+d_ny_n+e_ny_{n+1}=r_n
\end{equation}
의 해이고 그것이 큰 $n$에서 동차 점화식의 지배해보다 천천히 자라면 앞으로 점화는 같은 이유로 실패한다. 앞으로 계산하는 동안 들어온 반올림이 동차 해의 지배 성분을 키우기 때문이다. Olver\cite{olver1967difference}는 이 계산을 경계값 문제로 바꾸었다. 처음 값 $y_0$와 충분히 큰 $N$의 끝값 $y_N$(근삿값이나 $0$)을 주고 \cref{eq:inhom}의 $n=1,\dots,N-1$을 연립방정식으로 푼다. 그 계수 행렬은 삼중대각이므로 Gauss 소거로 $O(N)$ 연산에 풀린다. 끝값의 오차는 \cref{thm:miller}와 같은 이유로 안쪽으로 오면서 지배해의 크기 비만큼 줄어든다. 끝값 $y_N$에 오차 $\delta$를 주면 해는 $y_0$에서 $0$이고 $y_N$에서 $\delta$인 동차 해만큼 바뀌는데 그 해는 $\delta\,(f_ng_0-g_nf_0)/(f_Ng_0-g_Nf_0)$이다. $f$가 최소해이면 $f_ng_0-g_nf_0\approx-g_nf_0$이고 분모도 $-g_Nf_0$ 정도이므로 이 값은 $n\ll N$에서 $\delta\,g_n/g_N$ 정도로 작다. Olver의 원래 알고리즘은 소거를 앞으로 진행하며 끝값의 영향을 추정해 그것이 원하는 정확도 아래가 되는 $N$에서 멈춘다. \cref{alg:olver}에는 같은 생각을 $N$을 두 배씩 늘리는 간단한 꼴로 적었다. 고정된 $N$이면 소거 한 번이다.

\begin{algorithm}[htbp]
\begin{algorithmic}[1]
\Require \cref{eq:inhom}의 계수, $y_0$, 원하는 차수 $n_{\max}$, 허용오차 $\tau$, 끝값의 근사 $\eta(N)$ (점근식이나 $0$)
\Ensure $y_1,\dots,y_{n_{\max}}$
\State $N\gets2n_{\max}$; $y^{\mathrm{old}}\gets$ 없음
\Loop
  \State $n=1,\dots,N-1$의 방정식을 세운다: 첫 식에서 $c_1y_0$을, 끝 식에서 $e_{N-1}\eta(N)$을 오른쪽으로 옮긴다
  \State $y\gets$ \cref{alg:dgtsl}로 푼 해 \Comment{$O(N)$}
  \State \textbf{if} $y^{\mathrm{old}}$가 있고 $\max_{n\le n_{\max}}|y_n-y^{\mathrm{old}}_n|\le\tau\max|y_n|$ \textbf{then} \Return $y_{1..n_{\max}}$
  \State $y^{\mathrm{old}}\gets y$; $N\gets2N$
\EndLoop
\end{algorithmic}
\caption{끝값을 멀리 두는 경계값 풀이 (Olver의 방법을 단순화한 꼴)}
\label{alg:olver}
\end{algorithm}

\cref{alg:olver}의 정당성은 \cref{prop:discrete-green}의 거꾸로 가는 꼴에서 나온다. 끝값의 오차가 $g_n/g_N$으로 줄어들므로 $N$을 두 배로 했을 때 앞쪽 값이 거의 바뀌지 않으면 그 값들은 이미 끝값의 영향을 벗어났다. 비용은 $N$의 합, 곧 마지막 $N$의 두 배 정도이다.

\begin{algorithm}[htbp]
\begin{algorithmic}[1]
\Require 삼중대각 계(아래 대각 $c_{2..n}$, 대각 $d$, 위 대각 $e_{1..n-1}$), 오른쪽 $b$
\Ensure 해 $b$ (제자리)
\For{$k=1,\dots,n-2$}
  \If{$|d_k|<|c_{k+1}|$} \Comment{부분 피벗: 행 $k$와 $k+1$을 맞바꾼다}
    \State $t\gets-d_k/c_{k+1}$; $(d_k,d_{k+1})\gets(c_{k+1},\ \mathrm{fma}(t,d_{k+1},e_k))$
    \State $c_{k+1}\gets e_{k+1}$; $e_k\gets$ 옛 $d_{k+1}$; $e_{k+1}\gets te_{k+1}$ \Comment{$c_{k+1}$이 둘째 위 대각을 맡는다}
    \State $(b_k,b_{k+1})\gets(b_{k+1},\ \mathrm{fma}(t,b_{k+1},b_k))$
  \Else
    \State $t\gets-c_{k+1}/d_k$; $d_{k+1}\gets\mathrm{fma}(t,e_k,d_{k+1})$; $b_{k+1}\gets\mathrm{fma}(t,b_k,b_{k+1})$; $c_{k+1}\gets0$
  \EndIf
\EndFor
\State $b_n\gets b_n/d_n$; $b_{n-1}\gets(b_{n-1}-e_{n-1}b_n)/d_{n-1}$
\For{$k=n-2$ \textbf{downto} $1$} $b_k\gets(b_k-e_kb_{k+1}-c_{k+1}b_{k+2})/d_k$
\EndFor
\end{algorithmic}
\caption{부분 피벗을 쓰는 삼중대각 풀이 (LINPACK \code{dgtsl}의 scipy C 이식; stardust \code{csrc/integrate/quadpack.hpp})}
\label{alg:dgtsl}
\end{algorithm}

\cref{alg:dgtsl}은 LINPACK의 \code{dgtsl}이다. 부분 피벗으로 행을 맞바꾸면 아래 행의 첫 원소가 피벗이 되므로 위 대각이 하나 더 생긴다. 그 자리를 비어 있는 $c_{k+1}$에 두어 저장 공간을 늘리지 않는다. 비용은 나눗셈 $2n$번과 \code{fma} $4n$번 남짓이다. 피벗을 쓰므로 대각 우세가 아닌 계에서도 Gauss 소거의 표준적인 안정성 결과(48장)가 그대로 적용된다. stardust의 이 함수는 scipy의 C 이식을 따라 소거를 $n-2$행에서 멈추고 마지막 행은 아래 대각이 $0$인 것처럼 푼다. LINPACK 원본은 $n-1$행까지 소거하고 $|d_k|=|c_{k+1}|$인 동률에서도 행을 맞바꾸지만 scipy의 꼴은 엄격히 작을 때만 맞바꾼다. stardust는 scipy와 같은 비트를 내는 쪽을 택했다. 아래 측정에서 이 꼴의 경계값 풀이는 $p\in(2,24]$ 전체에서 모멘트의 크기에 대해 $1.4\times10^{-15}$ 안이었다(\cref{fig:moments-scan}).

\subsection{Chebyshev Moments in QUADPACK}

QUADPACK\cite{piessens1983quadpack}은 진동 적분 $\int_a^bf(x)\cos(\omega x)\,dx$를 계산할 때(stardust에서는 \code{integrate.quad}에 \code{weight='cos'}를 준 호출) 구간이 진동을 충분히 담으면 $f$를 $24$차 Chebyshev 급수로 근사하고 그 급수의 항마다 진동 무게를 곱한 적분을 정확히 더한다(\code{qc25f}). 이때 수정된 Chebyshev 모멘트 $\int_{-1}^1T_k(x)\cos(px)\,dx$($p=\omega h$, $h$는 구간의 반 길이)가 필요하다. 이 모멘트들은 $k$에 대한 비동차 삼항 점화식(같은 홀짝의 차수끼리)을 만족하고 처음 몇 개는 닫힌 꼴이 있다. 점화식은 Chebyshev 다항식의 미분방정식 $(1-x^2)T_k''-xT_k'+k^2T_k=0$과 $T_k'$의 관계 $2T_k=T'_{k+1}/(k+1)-T'_{k-1}/(k-1)$을 써서 $\int T_k\cos(px)\,dx$를 두 번 부분적분하면 나온다. 그 결과는 $\mu_{k-2}$, $\mu_k$, $\mu_{k+2}$와 끝점에서 나오는 $\sin p$, $p\cos p$의 항을 잇는 식이다(Piessens와 Branders, \cite{piessens1983quadpack}). stardust의 코드는 QUADPACK의 식을 항 하나하나 그대로 옮겼다. 동차 부분의 일반적인 해는 차수가 $|p|$를 넘어가면 계승처럼 빠르게 자란다($p=10$에서 차수 $14$--$24$ 구간은 차수를 $8$ 올릴 때마다 $10^4$--$10^6$배씩, 더 큰 차수에서는 더 빨리). 원하는 모멘트는 그 영역에서 천천히 줄어드므로 앞으로 점화는 차수가 $|p|$ 아래인 동안만 안정하다.

QUADPACK은 이 경계를 기준으로 두 경우를 나눈다. 필요한 차수는 $24$까지이므로 $|p|>24$이면 앞으로 점화가 안정하고 그대로 쓴다. $|p|\le24$이면 $25$개의 방정식을 경계값 문제로 푼다. 필요한 것은 차수 $24$까지인데 미지수는 차수 $6$부터 $54$까지이다. 끝을 필요한 차수의 두 배 너머에 두었으므로 \cref{alg:olver}의 원리대로 끝값의 오차는 안쪽에서 사라진다. 끝값은 큰 $k$에서의 점근 전개를 $1/k^2$의 네 항까지 계산한 것(코드의 \code{asap})이고 \cref{alg:dgtsl}로 푼다. 계산한 모멘트는 반복 이분의 같은 깊이에서 다시 쓰도록 표(\code{chebmo})에 저장해 둔다. stardust는 이 경로를 원본 그대로 옮겼다(\code{qc25f}, \code{dgtsl}). \cref{fig:moments}에 $p=10$과 $p=30$에서 잰 $\cos$ 모멘트의 오차를 그렸다. 참값은 같은 점화식을 십진 산술 $100$자리로 앞으로 돌려 얻었다. 그 정밀도에서는 증폭되어도 수십 자리가 남는다. $p=10$에서 binary64 앞으로 점화의 오차는 차수 $14$에서 $2.2\times10^{-14}$, $24$에서 $9.0\times10^{-8}$으로 자란다. 경계값 풀이(\code{qc25f}와 \code{dgtsl}을 옮긴 코드)는 모든 차수에서 $5.2\times10^{-16}$ 아래이다. $p=30$에서는 앞으로 점화가 차수 $24$까지 $2.4\times10^{-14}$ 아래로 머문다. 따로 빌드한 stardust의 \code{integrate.quad(exp, 0, 2, weight='cos', wvar=w)}를 $w=10,20,40,60,200$에서 정확한 값과 견주면 상대오차가 모두 $6.3\times10^{-16}$ 아래였다.

\begin{figure}[htbp]
\centering
\begin{tikzpicture}
\begin{semilogyaxis}[sdplot,xlabel={Chebyshev 차수 $k$},ylabel={모멘트의 상대오차},xmin=0,xmax=25,ymin=1e-17,ymax=1e-6,ytick={1e-16,1e-14,1e-12,1e-10,1e-8,1e-6}]
\addplot[black,mark=*,mark size=1.2pt] table[x=deg,y=fwd,restrict expr to domain={\thisrow{p}}{9:11}]{figures/06/moments.dat};
\addlegendentry{$p=10$, 앞으로}
\addplot[black,dashed,mark=square*,mark size=1.2pt] table[x=deg,y=bvp,restrict expr to domain={\thisrow{p}}{9:11}]{figures/06/moments.dat};
\addlegendentry{$p=10$, 경계값}
\addplot[black!50,mark=triangle*,mark size=1.4pt] table[x=deg,y=fwd,restrict expr to domain={\thisrow{p}}{29:31}]{figures/06/moments.dat};
\addlegendentry{$p=30$, 앞으로}
\end{semilogyaxis}
\end{tikzpicture}
\caption{$\int_{-1}^1T_k(x)\cos(px)\,dx$의 상대오차(\code{figures/06/moments.py}).}
\label{fig:moments}
\end{figure}

\cref{fig:moments-scan}에서는 같은 비교를 $p\in(2,60]$에서 훑었다. 모멘트 하나가 $0$ 근처를 지날 때는 그 상대오차에 의미가 없으므로 오차는 차수 $0$--$24$의 모멘트 가운데 가장 큰 것의 크기에 대해 쟀다. $p\le24$에서 앞으로 점화는 $p$가 작을수록 나빠져 $p=2.25$에서 $10^6$을 넘고 경계값 풀이는 어디서나 $1.4\times10^{-15}$ 아래이다. $p>24$에서 앞으로 점화는 최대 $1.1\times10^{-14}$이다.

\begin{figure}[htbp]
\centering
\begin{tikzpicture}
\begin{semilogyaxis}[sdplot,xlabel={$p$},ylabel={최대 오차/모멘트의 크기},xmin=2,xmax=60,ymin=1e-17,ymax=1e8,ytick={1e-16,1e-12,1e-8,1e-4,1,1e4,1e8},unbounded coords=jump]
\addplot[black,mark=*,mark size=0.9pt] table[x=p,y=fwd]{figures/06/moments_scan.dat};
\addlegendentry{앞으로 점화}
\addplot[black,dashed,mark=square*,mark size=0.9pt] table[x=p,y=bvp]{figures/06/moments_scan.dat};
\addlegendentry{경계값 ($p\le24$)}
\addplot[black!45,thin,forget plot] coordinates {(24,1e-17) (24,1e8)};
\end{semilogyaxis}
\end{tikzpicture}
\caption{$\cos$ 모멘트(차수 $0$--$24$)의 최대 오차를 $p$에 따라 그렸다. 세로선 $p=24$가 QUADPACK의 경계이다(\code{figures/06/moments\_scan.py}).}
\label{fig:moments-scan}
\end{figure}


\section{Recurrences in Other Parameters}
\label{sec:other-recurrences}

방향의 규칙은 차수에 대한 점화식에만 해당하지 않는다. 아래는 stardust에서 다른 매개변수에 대한 점화식을 쓰고 저마다 방향을 판정하는 예이다.

\subsection{Kummer's Function in Its First Parameter}

합류 초기하 함수 $M(a,b,x)={}_1F_1(a;b;x)$(93장)는 $a\le-2$, $b>0$에서 거듭제곱 급수의 항들이 $e^{2\sqrt{|a|x}}$처럼 자랐다가 상쇄한다. stardust는 그곳에서 $a$에 대한 점화식(\cite{dlmf} 13.3.1)
\[
(b-a)M(a-1)=(b-2a-x)M(a)+aM(a+1)
\]
을 아래로 돌린다(\code{csrc/special/ops/hyper.hpp}의 \code{m\_a\_recurrence}). 어느 방향이 안정한지는 $k=-a$가 $k_0=(x-b)^2/(4x)$보다 작은지 큰지로 정해진다($x>b$). $-a<k_0$인 구간에서는 $M(a-k)$가 아래로 갈수록 사라지는 해, 곧 그 방향의 최소해이므로 점화가 그것을 잃는다. $k_0$부터는 그렇지 않다. 두 해의 특성비를 곱하면 $k/(b+k)<1$이고 $M$은 두 해 가운데 큰 쪽을 따라가다가 그 뒤 둘 다 진동한다. 그래서 시작점을 $k_0$ 너머에 둔다. 다항식($a=-n$)이면 $M(0)=1$, $M(-1)=1-x/b$에서 바로 시작하고 그 밖이면 $a+t$와 $a+t-1$($a+t$가 $-k_0$ 바로 아래)에서 보정 합의 급수로 두 값을 만든다.

\begin{algorithm}[htbp]
\begin{algorithmic}[1]
\Require $a=a_h+a_l\le-2$ (double-double), $b>0$, $x$
\Ensure $M(a,b,x)$, 또는 ``이 방법으로는 안 됨''
\State 시작 $\tau$를 고르고 $(M(\tau),M(\tau-1))$을 double-double로 만든다 \Comment{다항식이면 $\tau=0$}
\State $(D_{\mathrm{prev}},D)\gets(0,1)$ \Comment{같은 점화식의 둘째 해, binary64: 해의 크기를 잰다}
\For{$k=1,\dots,\tau-a_h-1$}
  \State $a_k\gets\tau-k$ (정확); 분모 $b-a_k-a_l$을 double-double로
  \State $M\gets((b-x-2a_k-2a_l)M+(a_k+a_l)M_{\mathrm{prev}})/(b-a_k-a_l)$ \Comment{double-double}
  \State $D\gets((b-x-2a_k)D+a_kD_{\mathrm{prev}})/(b-a_k)$ \Comment{binary64}
  \State 두 해를 각각 $2^{\pm600}$으로 척도 맞추고 지수를 센다
\EndFor
\State $\varepsilon\gets(\text{걸음 수})\cdot2^{-102}+(\text{시작값의 오차})$
\State \textbf{if} $\varepsilon\cdot\max(|D|,|M_{\mathrm{prev}}|)>2^{-54}|M|$ \textbf{then} \Return 안 됨 \Comment{결과가 해의 크기보다 너무 작다}
\State \Return $M$
\end{algorithmic}
\caption{$a$에 대한 점화로 하는 $M(a,b,x)$ (stardust \code{hyper.hpp}의 \code{m\_a\_recurrence})}
\label{alg:m-a}
\end{algorithm}

\cref{alg:m-a}는 6행에서 둘째 해 $D$를 함께 돌린다. 걸음마다의 반올림은 \cref{prop:discrete-green}에 따라 해들의 크기에 비례하는 오차를 남긴다. 그러므로 잃는 정밀도는 결과 $M$이 해들의 크기보다 얼마나 작은지로 정해지고 $M$이 영점 근처에 있을 때 크다. $(0,1)$에서 시작한 $D$를 binary64로 함께 돌리면 해의 크기를 싸게 잴 수 있고 10행이 그 비로 남은 비트를 판정한다. 남은 비트가 $53$보다 적으면 다른 방법(Boost, 점근 전개)으로 넘긴다. $b<0$에서는 이 방법을 쓰지 않는다. 둘째 해가 $k=-b/2$까지 줄었다가 그 뒤 자라므로 중간에 생긴 반올림이 증폭되는데 $D$가 그것을 보여 주지 못하기 때문이다. 코드의 주석에 남은 비교로는 Boost의 binary64 점화가 $a=-10^5$에서 $2\times10^{-9}$ 틀렸고 $b<0$에서는 몇 자릿수씩 틀렸다.

\subsection{Discrete Orthogonal Polynomials}

Savitzky--Golay 필터(\code{smooth.savgol\_coeffs}, 130장)의 계수는 창의 점 $s_i=i-(N-1)/2$($i<N$) 위의 이산 직교 다항식(Gram 다항식) $q_j$로 쓰인다. 차수 $d$의 최소제곱 다항식의 $\nu$차 도함수를 위치 $s_p$에서 내는 계수는 $h_i=\delta^{-\nu}\sum_{j=\nu}^dq_j(s_i)q_j^{(\nu)}(s_p)$이다. $q_j$는 차수에 대한 삼항 점화식 $s\,q_j=\alpha_{j+1}q_{j+1}+\alpha_jq_{j-1}$을 만족하므로 $q_j(s_i)$를 그 점화식으로 만드는 것이 자연스러워 보인다. 그런데 창의 양 끝은 높은 차수의 다항식에서 지수적으로 작은 영역에 있다. 차수가 $N/4$쯤을 넘으면 끝점 $s_0$이 그 다항식들의 전이점 바깥에 놓인다. 그곳에서 차수 방향의 점화식은 최소해를 앞으로 계산하는 셈이 되어 \cref{prop:forward-error}대로 망가진다.

\cref{fig:gram}에 $N=51$에서 단항식 꼴로 정규화한 Gram 다항식 $p_j$를 차수 점화식으로 binary64 계산한 상대오차를 그렸다. 참값은 정확한 유리수이다. 창의 가운데에서는 오차가 반올림 수준에 머문다. 끝점에서는 차수 $j=32$에서 $10^{-8}$을 넘고 $j=45$에서는 $3.2\times10^3$이다. stardust는 같은 값을 점의 방향 $i$로 점화해 얻는다(Mukundan의 방법). 창의 첫 점과 그 이웃의 닫힌 꼴에서 시작해 가운데까지 가고 나머지는 대칭 $q_j(s_{N-1-i})=(-1)^jq_j(s_i)$로 채운다. 이 방향에서는 값이 작은 끝에서 큰 가운데로 자라므로 지배해를 앞으로 계산하는 셈이다. 계산은 double-double로 한다. 따로 빌드한 stardust의 \code{savgol\_coeffs}$(51,d)$를 정확한 유리수 계수와 견주면 $d=10,20,30,40$에서 최대 오차가 계수의 최댓값에 대해 $5.4\times10^{-17}$, $8.5\times10^{-17}$, $3.1\times10^{-17}$, $3.6\times10^{-17}$이었다. 같은 합을 차수 점화식의 binary64 값으로 만들면 $2.3\times10^{-16}$, $2.9\times10^{-16}$, $1.8\times10^{-14}$, $6.1\times10^{-11}$이었다(\code{figures/06/gram.py}).

\begin{figure}[htbp]
\centering
\begin{tikzpicture}
\begin{semilogyaxis}[sdplot,xlabel={차수 $j$},ylabel={$p_j(s)$의 상대오차},xmin=0,xmax=46,ymin=1e-17,ymax=1e4,ytick={1e-16,1e-12,1e-8,1e-4,1,1e4}]
\addplot[black,mark=*,mark size=1.1pt] table[x=j,y=end]{figures/06/gram.dat};
\addlegendentry{끝점 $s_0$}
\addplot[black!50,mark=square*,mark size=1.1pt] table[x=j,y=centre]{figures/06/gram.dat};
\addlegendentry{가운데 $s_{25}$}
\addplot[black!45,thin,forget plot] coordinates {(12.75,1e-17) (12.75,1e4)};
\end{semilogyaxis}
\end{tikzpicture}
\caption{$N=51$ 창의 Gram 다항식을 차수 점화식으로 binary64 계산한 상대오차(\code{figures/06/gram.py}). 세로선은 $N/4$.}
\label{fig:gram}
\end{figure}

\subsection{First-Order Recurrences}

일계 점화식에도 방향이 있다. 비중심 카이제곱, $F$, $t$ 분포의 확률은 Poisson 가중치를 단 불완전 함수들의 합 $\sum_jw_jG(s_0+j)$이다(\code{csrc/special/ops/noncentral.hpp}). 이웃한 $G$는 $G(s+1)=G(s)\mp T(s)$로 이어지고 $T(s)=|G(s)-G(s+1)|\ge0$은 \cref{alg:gamma-prefix}의 앞 인자 꼴로 직접 계산된다. 두 방향은 오차가 쌓이는 모양이 다르다.

\begin{proposition}
\label{prop:first-order}
$y_{k+1}=y_k+t_k$, $t_k\ge0$, $y_0\ge0$을 binary64로 계산하면 $\hat y_n$의 상대오차는 $\gamma_n$ 정도($t_k$ 자신의 상대오차 $\eta$가 있으면 $\gamma_n+\eta$) 이하이다. 거꾸로 $y_k=y_{k+1}-t_k$로 계산하면 상대오차는 $\sum_{j\ge k}t_j/y_k$배로 커질 수 있다.
\end{proposition}

\begin{proof}
앞 방향은 음이 아닌 수들의 합이므로 합의 표준적인 오차 한계(11장)에서 상대오차가 $\gamma_n$ 이하이다. 거꾸로 가면 $y_k=y_n-\sum_{j=k}^{n-1}t_j$가 되어 큰 수들의 차로 작은 수를 만들므로 각 항의 절대오차 $u\,t_j$가 $y_k$에 대해 $u\sum t_j/y_k$배로 보인다.
\end{proof}

그래서 stardust는 $T$를 더하는 방향으로만 걷는다. 걸음은 반대쪽의 먼 끝에서 시작해 가장 큰 항들을 지나 항이 사라지는 곳까지 간다. 시작점에서만 $w$, $G$, $T$를 직접 계산하고 나머지는 비로 이어 간다. $w$와 $T$의 비는 ``double과 오차''를 함께 들고 다니는 산술(걸음마다 약 $2^{-100}$)로 곱해 반올림이 표류하지 않게 한다. 코드의 주석에 따르면 scipy가 부르는 Boost의 판은 꼬리에서 $0$을 내거나 $8.8\times10^{15}$ ulp까지 틀렸다.

이산 분포의 엔트로피와 꼬리 합(\code{csrc/stats/discrete.hpp})도 같은 꼴이다. 확률 질량의 비 $\mathrm{pmf}(j+1)/\mathrm{pmf}(j)$는 정수와 모양 모수의 곱과 몫 네 개로 쓰이므로 각 인자를 double-double로 정확히 만들고 비 점화 $t_{j+1}=t_j\cdot(a_1a_2)/(b_1b_2)$를 double-double로 돌린다. 걸음마다 오차가 $2^{-104}$ 정도이므로 $2^{25}$걸음을 가도 binary64가 보는 범위에서 잃는 것이 없다. 로그 질량을 여덟 걸음마다 다시 계산해 맞출 필요도 없다. 걸음은 최빈값에서 바깥으로 간다. 이 방향에서 항이 줄어드므로 합이 앞 방향의 합이 된다.

\subsection{Orthogonal Polynomials for Gauss Quadrature}

\code{integrate.fixed\_quad}(76장)의 Gauss--Legendre 마디는 $n\le100$이면 Arb로 만든 표이고 그보다 크면 Newton 방법으로 $P_n(x)=0$을 푼다(\code{csrc/integrate/gauss.cpp}). $P_n$과 $P_n'$은 점화식 $(j+1)P_{j+1}=(2j+1)xP_j-jP_{j-1}$과 $P_n'=n(xP_n-P_{n-1})/(x^2-1)$로 계산한다. $|x|<1$에서 이 점화식의 두 해($P_n$과 둘째 종류 $Q_n$)는 비슷한 크기로 진동하므로 점화는 중립이고 반올림만 쌓인다. stardust는 binary64로 근사 마디를 찾은 뒤 double-double로 두 번 더 Newton 걸음을 하고 마디와 무게를 올바르게 반올림한다. 마디마다 따로 계산하므로 스레드 수와 무관한 비트가 나온다.

\section{Clenshaw's Algorithm}
\label{sec:clenshaw}

\subsection{Summing Series of Recurrent Functions}

$\phi_k$가 삼항 점화식 $\phi_{k+1}=\alpha_k(x)\phi_k+\beta_k(x)\phi_{k-1}$을 만족할 때 급수 $S=\sum_{k=0}^nc_k\phi_k(x)$를 계산한다고 하자. 직교 다항식 급수가 대표적이다. $\phi_k$를 하나씩 만들어 곱해 더할 수도 있지만 Clenshaw\cite{clenshaw1955}는 점화식을 거꾸로 돌려 계수를 접는 방법을 내놓았다. Chebyshev 다항식 $T_{k+1}=2tT_k-T_{k-1}$($T_0=1$, $T_1=t$)이면 다음과 같다.

\begin{proposition}
\label{prop:clenshaw}
$b_{n+1}=b_{n+2}=0$, $b_k=c_k+2tb_{k+1}-b_{k+2}$($k=n,\dots,1$)이라 하면 $\sum_{k=0}^nc_kT_k(t)=c_0+tb_1-b_2$이다.
\end{proposition}

\begin{proof}
$k\ge1$에서 $c_k=b_k-2tb_{k+1}+b_{k+2}$를 합에 넣고 $b_j$마다 계수를 모으면 $j\ge3$의 $b_j$에는 $T_j-2tT_{j-1}+T_{j-2}=0$이 곱해진다. 남는 것은 $b_1T_1+b_2(T_2-2tT_1)=tb_1-b_2$이다($T_2-2t^2=-1$). 여기에 $c_0T_0=c_0$을 더한다.
\end{proof}

비용은 곱셈 $n$번($2t$는 미리 계산)과 덧셈 $2n$번이고 $T_k$를 따로 만들지 않는다. stardust에는 Clenshaw 합이 두 곳에 있다. 하나는 Cephes의 \code{chbevl}(\code{csrc/special/ops/bessel01.hpp})로 $e^{-x}I_0(x)$, $e^{-x}I_1(x)$, $e^xK_0(x)$, $e^xK_1(x)$를 구간마다 Chebyshev 급수로 근사한 계수를 더한다. Cephes의 관례에 따라 인자를 $T_k(t/2)$ 꼴로 받고 첫 계수를 마지막에 반으로 나눈다. 다른 하나는 다변량 $t$ 분포의 준몬테카를로 적분(\code{csrc/stats/mvcdf.hpp}의 \code{Cheb})이다. $\chi$ 변수의 역함수 $r(w)$를 매번 \code{gammaincinv}로 구하는 대신 조각마다 $16$개 마디의 Chebyshev 보간을 만들어 두고 이 합으로 계산한다. 코드의 주석에 따르면 이 표는 값 하나에 $15$ ns로 \code{gammaincinv}의 $260$ ns를 대신하고 $5$차원 문제 하나의 시간을 $183$ ms에서 $37$ ms로 줄였다.

\begin{algorithm}[htbp]
\begin{algorithmic}[1]
\Require 계수 $c_0,\dots,c_n$, $t\in[-1,1]$
\Ensure $\sum_{k=0}^nc_kT_k(t)$ (\code{mvcdf}의 꼴: $c_0/2$를 쓰면 마지막 줄의 $c_0$을 $c_0/2$로)
\State $b_1\gets0$; $b_2\gets0$
\For{$k=n$ \textbf{downto} $1$}
  \State $b_0\gets2tb_1-b_2+c_k$; $b_2\gets b_1$; $b_1\gets b_0$
\EndFor
\State \Return $c_0+(tb_1-b_2)$ \Comment{상수항은 마지막에 더한다}
\end{algorithmic}
\caption{Clenshaw 합 (stardust \code{mvcdf.hpp}의 \code{Cheb::eval}, \code{bessel01.hpp}의 \code{chbevl})}
\label{alg:clenshaw}
\end{algorithm}

\subsection{Rounding Errors}

\begin{proposition}
\label{prop:clenshaw-error}
\cref{alg:clenshaw}을 binary64로 계산한 결과 $\hat S$는 계수 $c_k$를 $c_k+\eta_k$로 바꾼 급수의 정확한 값과 마지막 줄의 반올림 몇 번 안에서 같다. 여기서 $|\eta_k|\le\gamma_3(|c_k|+2|t||\hat b_{k+1}|+|\hat b_{k+2}|)$이다. 따라서 $|\hat S-S|\lesssim\gamma_3\sum_k(|c_k|+3|\hat b_k|)$(마지막 줄의 반올림 포함)이고 $|b_k|\le\sum_{j\ge k}(j-k+1)|c_j|$이다.
\end{proposition}

\begin{proof}
$k$번째 줄의 계산은 $\hat b_k=(c_k+2t\hat b_{k+1}-\hat b_{k+2})(1+\theta)$ 꼴이고 표준적인 한계(11장)에서 $\hat b_k=c_k+2t\hat b_{k+1}-\hat b_{k+2}+\eta_k$로 쓸 수 있다. $\eta_k$를 $c_k$에 넣으면 계산된 $\hat b$들은 바뀐 계수 $c_k+\eta_k$에 대한 정확한 점화가 된다. \cref{prop:clenshaw}에서 그 결과는 $\sum(c_k+\eta_k)T_k$이고 $|T_k|\le1$이므로 오차는 $\sum|\eta_k|$ 이하이다. $b_k$는 $c$에 대해 선형이고 $c_j$가 $b_k$에 주는 기여는 $c_jU_{j-k}(t)$이다($U_m$은 둘째 종류 Chebyshev 다항식으로 같은 점화식을 $U_0=1$, $U_1=2t$로 시작한 것). $|U_m(t)|\le m+1$이므로 마지막 부등식이 나온다.
\end{proof}

계수가 빨리 줄면 $\sum|b_k|$가 $\sum|c_k|$의 작은 배수이고 오차는 $u\sum|c_k|$ 정도이다. $|U_m|$의 한계 $m+1$은 $t=\pm1$에서 이루어지므로 계수가 천천히 줄면 끝점 근처에서 오차가 차수의 제곱에 비례해 커질 수 있다. 이 정도면 충분히 안정하다. 같은 다항식을 단항식 $\sum m_kt^k$로 바꾸어 Horner로 계산해 보면 차이가 크다. 단항식 계수는 $T_k$의 선행 계수가 $2^{k-1}$이므로 커지고 Horner의 오차는 $u\sum|m_k||t|^k$ 정도이다. \cref{fig:clenshaw}에 계수가 $c_k=\pm0.85^k$($k\le40$, 부호는 무작위)인 다항식을 두 방법으로 계산한 절대오차를 그렸다. 참값은 정확한 유리수로 계산했다. $\sum|c_k|=6.7$인 이 다항식의 단항식 계수는 최대 $2.3\times10^{11}$이다. Clenshaw의 오차는 $[-1,1]$ 전체에서 $1.2\times10^{-15}$ 아래이고 Horner는 끝점에서 $5.4\times10^{-5}$이다. 따로 빌드한 stardust의 \code{special.i0e}는 \code{chbevl}의 두 계수 표가 갈리는 $x\le8$과 $x>8$에서 십진 산술의 급수와 견주어 최대 상대오차가 각각 $4.4\times10^{-16}$, $4.7\times10^{-16}$이었다.

\begin{figure}[htbp]
\centering
\begin{tikzpicture}
\begin{semilogyaxis}[sdplot,xlabel={$t$},ylabel={절대오차},xmin=-1,xmax=1,ymin=1e-18,ymax=1e-3,ytick={1e-17,1e-14,1e-11,1e-8,1e-5}]
\addplot[black,mark=*,mark size=0.8pt] table[x=t,y=clenshaw]{figures/06/clenshaw.dat};
\addlegendentry{Clenshaw}
\addplot[black!50,mark=square*,mark size=0.8pt] table[x=t,y=horner]{figures/06/clenshaw.dat};
\addlegendentry{단항식 Horner}
\end{semilogyaxis}
\end{tikzpicture}
\caption{$40$차 Chebyshev 급수를 Clenshaw 합과 단항식의 Horner 계산으로 구한 절대오차(\code{figures/06/clenshaw.py}; $10^{-18}$은 정확히 맞은 점).}
\label{fig:clenshaw}
\end{figure}


\section{Testing Recurrence Code}
\label{sec:rec-testing}

점화식을 쓰는 코드를 시험하는 길은 셋이다. 가장 직접적인 것은 높은 정밀도의 참값이다. 이 장에서는 mpmath 없이 십진 산술로 급수(\cref{eq:jv-series})나 같은 점화식을 $90$--$100$자리로 돌려 참값을 만들었다(\code{figures/06/dref.py}). 같은 점화식을 높은 정밀도로 돌리는 것은 불안정한 방향이어도 증폭된 오차가 그 정밀도 안에 머무는 동안에는 믿을 만하다. 다음으로 Wronskian 같은 항등식을 쓴다. 참값 없이 계산값끼리의 일관성을 보므로 넓은 범위를 싸게 훑을 수 있다. 다만 항등식에 쓰는 값들이 같은 방법에서 나왔다면 공통의 오차를 보지 못한다. 항등식에 넣는 인자가 정말 그 항등식의 인자인지도 확인해야 한다. 이 장의 측정에서 차수 $511.3$과 $512.3$의 Wronskian 잔차가 $1.3\times10^{-13}$으로 튀었는데 라이브러리가 아니라 측정의 잘못이었다. 두 수는 다른 이진 지수 구간에 있어 소수 부분의 반올림이 다르므로 double $512.3$과 $511.3$의 차는 정확히 $1$이 아니다. 차수의 차 $\delta\nu\approx5\times10^{-14}$에 $\partial_\nu\ln Y_\nu\approx\ln(2\nu/x)$를 곱한 만큼이 잔차에 나타났다. 차수를 $2^{-10}$의 배수로 잡자 잔차가 $10^{-15}$ 아래로 내려갔다. \cref{fig:wronsk}에는 그렇게 잡은 차수 $v\in[0.25,1000]$에서 잰 \code{jv}, \code{yv}의 Wronskian 잔차를 그렸다. 빠른 경로($v<500$)와 일반 경로 모두 $4\times10^{-15}$ 아래이고 $x=10$의 큰 차수는 $J$가 아래로 넘쳐 잔차를 잴 수 없다.

\begin{figure}[htbp]
\centering
\begin{tikzpicture}
\begin{semilogyaxis}[sdplot,xlabel={차수 $v$},ylabel={Wronskian 잔차},xmin=0,xmax=1000,ymin=1e-17,ymax=1e-13,ytick={1e-17,1e-16,1e-15,1e-14,1e-13},unbounded coords=jump]
\addplot[black,mark=*,mark size=0.9pt] table[x=nu,y=x10]{figures/06/wronsk.dat};
\addlegendentry{$x=10$}
\addplot[black,dashed,mark=square*,mark size=0.9pt] table[x=nu,y=x100]{figures/06/wronsk.dat};
\addlegendentry{$x=100$}
\addplot[black!50,mark=triangle*,mark size=1.1pt] table[x=nu,y=x300]{figures/06/wronsk.dat};
\addlegendentry{$x=300$}
\addplot[black!45,thin,forget plot] coordinates {(500,1e-17) (500,1e-13)};
\end{semilogyaxis}
\end{tikzpicture}
\caption{따로 빌드한 stardust의 \code{jv}, \code{yv}로 잰 Wronskian 잔차 $|(\pi x/2)(J_{v+1}Y_v-J_vY_{v+1})-1|$을 곱의 크기로 나눈 값(\code{figures/06/wronsk.py}; $10^{-17}$은 그 이하). 세로선은 빠른 경로의 상한 $500$.}
\label{fig:wronsk}
\end{figure} stardust의 \code{jy\_temme}는 Wronskian으로 정규화하므로 그 경로의 $J$와 $Y$는 Wronskian을 거의 저절로 만족한다. 그 때문에 이 장에서는 \code{jv}를 십진 급수와도 따로 견주었다. 마지막으로 경로의 경계를 살핀다. stardust의 Bessel 함수는 영역마다 다른 방법(급수, Hankel, 점화, Miller, Temme, Steed, Debye)을 쓰므로 영역의 경계 양쪽에서 값이 이어지는지를 본다. 라이브러리의 정확도 시험은 구간 산술(Arb)로 만든 참값 파일과 견준다(\code{tests/special/test\_reference.py}). \cref{tab:rec-tests}에 이 장의 측정을 모았다. 모든 라이브러리 수치는 공개 저장소를 따로 받아 빌드한 stardust로 쟀다.

\begin{table}[htbp]
\begin{sdtable}{colspec={lX}, row{1}={font=\small\bfseries}}
\toprule
대상 & 기준과 결과 \\
\midrule
앞으로 점화의 $J_n(1)$ & 십진 급수; $n=10$에서 상대오차 $27$, $J_{20}$은 부호까지 틀림 \\
Miller ($N=30$) & 같은 참값; $n\le25$에서 $3.1\times10^{-16}$ 이하 \\
Miller의 시작 차수 & $x=10$; 여유 $24$(차수 $0.3$), $14$(차수 $20.3$)에서 반올림 수준, stardust는 $30.0$, $35.7$ \\
\code{special.jv} & 십진 급수; 포락선 기준 $9.0\times10^{-16}$, Wronskian 잔차 $1.2\times10^{-15}$ ($\nu\le40$) \\
\code{jv}, \code{yv} ($v\le1000$) & Wronskian 잔차; $4\times10^{-15}$ 이하, 차수 $511.3$의 가짜 잔차는 측정의 잘못 \\
긴 점화 & $400$걸음; binary64 $31$--$40u$, 빠른 경로 식 $32$--$43u$, double-double $1.3u$ 이하 \\
CF1 & 항 수가 $x>\nu$에서 $x$에 비례 \\
Neumann 합 & 조건수 $2.0$($\mu=0.3$, $x=10$)--$10.7$($\mu=-0.4$, $x=100$) \\
\code{special.airy} & Wronskian $\pi W-1$; $1.8\times10^{-15}$ \\
Chebyshev 모멘트 & 십진 점화; 앞으로($p=10$) $9\times10^{-8}$, 경계값은 $p\in(2,24]$에서 $1.4\times10^{-15}$ 이하 \\
\code{smooth.savgol\_coeffs} & 정확한 유리수; $d=40$에서 $3.6\times10^{-17}$ (차수 점화 binary64는 $6.1\times10^{-11}$) \\
\code{integrate.quad} (\code{weight='cos'}) & 정확한 적분; $6.3\times10^{-16}$ 이하 \\
Clenshaw 대 Horner & 정확한 유리수; $1.2\times10^{-15}$ 대 $5.4\times10^{-5}$ \\
\code{special.i0e} & 십진 급수; $4.7\times10^{-16}$ 이하 \\
\bottomrule
\end{sdtable}
\caption{이 장의 측정.}
\label{tab:rec-tests}
\end{table}


\section{APIs in This Chapter}
\label{sec:rec-apis}

\cref{tab:rec-apis}에 이 장의 커널을 쓰는 stardust API와 그 커널을 설명한 곳을 모았다. Bessel 함수의 영역 나누기와 다른 근사(Hankel, Debye, 유리 근사, Boost 경로)는 91--92장, \code{expm}은 57장, QUADPACK의 나머지 루틴은 76--77장에서 다룬다.

\begin{table}[htbp]
\begin{sdtable}{colspec={lXl}, row{1}={font=\small\bfseries}}
\toprule
API & 이 장에서 다룬 것 & 설명한 곳 \\
\midrule
\code{special.jv}, \code{jve} & 급수 영역, 위로 점화(중립), Miller와 Neumann 합 & \cref{alg:miller-bessel} \\
\code{special.jv}, \code{yv} (일반 경로) & CF1 뒤 패스와 부호, 아래로 점화, Wronskian, 위로 점화 & \cref{alg:cf1,alg:jy} \\
\code{special.yv}, \code{yve} (빠른 경로) & Temme 급수, CF2와 Miller의 $J$, Wronskian, 위로 점화 & \cref{alg:temme,alg:fast-y} \\
\code{special.iv}, \code{kv}, \code{ive}, \code{kve} & $K$ 위로 점화, $I$의 역방향 비, Wronskian & \cref{alg:ik} \\
\code{special.airy} & 음의 인자의 두 격자 Miller, Wronskian 검사 & \cref{sec:miller} \\
\code{special.hyp1f1} & $a$에 대한 점화와 둘째 해로 하는 판정 & \cref{alg:m-a} \\
\code{smooth.savgol\_coeffs}, \code{savgol\_filter} & Gram 다항식의 점 방향 점화 & \cref{sec:other-recurrences} \\
\code{stats}의 비중심 분포, 이산 분포의 엔트로피 & 더하는 방향의 일계 점화, 비 점화 & \cref{prop:first-order} \\
\code{integrate.fixed\_quad} & Legendre 점화와 double-double Newton & \cref{sec:other-recurrences} \\
\code{integrate.quad} (\code{weight='cos'}, \code{'sin'}) & Chebyshev 모멘트의 경계값 풀이 & \cref{alg:dgtsl} \\
\code{special.i0}, \code{i1}, \code{k0}, \code{k1}과 \code{*e} & Chebyshev 급수의 Clenshaw 합 & \cref{alg:clenshaw} \\
\code{stats.MultivariateT.cdf} & $\chi$ 역함수의 조각별 Chebyshev 표 & \cref{alg:clenshaw} \\
\code{linalg.expm} & 상수 계수 연립방정식의 해(이론) & \cref{prop:expm} \\
\bottomrule
\end{sdtable}
\caption{이 장의 커널을 쓰는 stardust API.}
\label{tab:rec-apis}
\end{table}

\section{Concluding Remarks}

이 장에서 본 결과는 모두 두 해가 다르게 자란다는 사실에서 나온다. 미분방정식에서는 Wronskian으로 두 해의 독립성을 재고 전이점을 경계로 진동과 지수적 분리가 나뉜다. 점화식에서는 Casorati 행렬식이 같은 일을 하고 구하려는 해가 최소해인지 지배해인지에 따라 계산 방향이 정해진다. 최소해는 거꾸로, 지배해는 앞으로 계산한다. 거꾸로 시작할 값이 없으면 아무 값에서 시작해 비율만 믿고 정규화 관계로 척도를 맞춘다(Miller). 그 비율은 연분수로도 얻는다(Pincherle). 원하는 해가 어느 쪽도 아니면 양 끝을 주고 연립방정식으로 푼다(Olver). 점화식으로 급수를 더할 때도 점화식을 거꾸로 돌린다(Clenshaw). 구현에서는 계수를 한 번만 반올림하고 차수를 정확히 만들며 긴 점화를 double-double로 하고 지수를 따로 들고 다녀야 정확도가 나왔다. 다음 장에서는 조합론과 생성 함수를 다룬다. 생성 함수의 계수는 다시 점화식을 만족하고 그 특이점의 위치가 계수의 증가 속도를 정한다.

\chapterreferences
